\chapter{Evaluation}
\label{ch:evaluation}

This chapter evaluates whether MERIT actually addresses the three gaps that were pointed out
earlier in Section~\ref{sec:relevance-to-merit}. 
Each of MERIT's pillars maps to a research question (RQ), which is derived from the literature
review earlier. Our studies are presented bottom-up (Ingestion $\rightarrow$ Fusion $\rightarrow$ 
Explainability $\rightarrow$ Human Alignment $\rightarrow$ Fairness and Scalability). Note that
this may disagree with the numerical prefixes of the studies in the repository, which follow the 
implementation order instead. Graphical representations of the studies are deferred to the appendix, 
and also cited inline.

\section{Evaluation Overview}
\label{sec:eval-overview}

\subsection{Evaluation Design and Research Questions}
\label{subsec:eval-design-rqs}

Our evaluation design is driven by the three formal research questions (RQ1--RQ3) established at the conclusion of our literature review (Section~\ref{sec:relevance-to-merit}). To briefly summarise, these questions evaluate MERIT's extensibility against rigid baselines, its multi-source fusion capabilities in the presence of conflicting candidate evidence, and the efficacy of its glass-box transparency for human decision calibration.

We map each RQ to prior work along with the associated studies and primary metrics to answer these RQs (Table~\ref{tab:eval-literature-traceability}).
We also report criteria that we use to determine the success of each RQ. Note that our success criteria are not
derived from literature, and are instead internal targets, since many results are directional due to the small sample size
in our evaluation. Both the Traditional ATS and Modern AI ATS baselines are simplified in-repository proxies
for the keyword-matching and AI-assisted ATS categories mentioned earlier in Section~\ref{sec:automated-screening}.
We do not implement a rule-based ATS baseline, the category is knockout-oriented and does not produce ranked, 
pool-wide scores comparable across candidates. Our studies instead require rankings 
(Spearman $\rho$), so we compare MERIT only against the two ranking-oriented CV-screening paradigms above. 
We do not believe this to be a significant threat to validity due to the similarities between Keyword-Matching and Rule-Based ATS systems.
For Studies~07--10, expert-consensus rankings were produced by the same four-recruiter panel described in Section~\ref{subsec:spearman-method},
who reviewed CV PDFs only,
following Lo et al.'s~\cite{lo2025aihiring} paradigm of comparing automated scores with HR judgement.
High Spearman $\rho$ measures agreement with that consensus, but cannot be generalised due to the sample size.

Our evaluation takes inspiration from software-engineering experimentation~\cite{wohlin_threats}
through construct isolation, where we use synthetic JSON personas to isolate scoring logic from 
the parsing noise of the ingestion engine.
We also take inspiration from Lo et al.'s~\cite{lo2025aihiring} evaluation paradigm, where they 
anonymise real resumes with an HR panel to evaluate the agreement between automated scores and HR judgement.
MERIT follows a related approach, where we create synthetic profiles based on real candidate data,
and then evaluate the agreement between the automated scores and HR judgement.
Threats to validity, including this design trade-off, are discussed in 
Section~\ref{sec:threats-to-validity}.


Against the internal criteria in Table~\ref{tab:eval-literature-traceability} (below), 
we do not treat every study as a pass. Study~14 meets the SUS target and Study~07 meets the directional $\rho \geq 0.7$ criterion for MERIT (Full).
However, failures are also documented: Studies~04--05 fall short of parsing $\geq 80\%$ and Spearman cohorts for Studies~08 and~10, 
fall below $\rho = 0.7$.
Outcomes by pillar are summarised in Section~\ref{sec:objectives-retrospective}.

\subsection{Evaluation Study Mapping}
\label{subsec:eval-codebase-traceability}

Each study maintains a one-to-one mapping with a prefixed subdirectory under \texttt{evaluation/} 
in the project repository~\cite{merit_repo} (e.g.\ Study~04 $\leftrightarrow$ \texttt{04-ir\_parser\_test/}).
Table~\ref{tab:eval-study-catalogue} details the purpose of each study. 
Note that Studies~02, ~03 and~06 are engineering supplements to either show proof of concept, or provide benchmarking data.

\begin{table}[H]
    \centering
    \caption[Evaluation study catalogue]{Mapping of each study to a specific research question (RQ) and purpose.}
    \label{tab:eval-study-catalogue}
    \footnotesize
    \renewcommand{\arraystretch}{1.28}
    \setlength{\tabcolsep}{5pt}
    \setlength{\extrarowheight}{1pt}
    \def\tabularxcolumn#1{m{#1}}%
    \begin{tabularx}{\textwidth}{@{}%
        >{\centering\arraybackslash}m{0.65cm}
        >{\centering\arraybackslash\hyphenpenalty=10000\exhyphenpenalty=10000}m{2.1cm}
        >{\centering\arraybackslash}X
        >{\centering\arraybackslash}m{0.48cm} @{}}
        \hline
        \textbf{No.} & \textbf{Name} & \textbf{Purpose} & \textbf{RQ} \\
        \hline
        \hyperref[subsec:study-01-comparative-ablation]{01} & Comparative ablation &
        Compare four engines on 10 personas and analyse fusion rank shifts. &
        2 \\
        \hyperref[subsec:studies-02-03-scalability]{02} & Runtime &
        Check runtime stays practical as $N$ grows (engineering supplement). &
        --- \\
        \hyperref[subsec:studies-02-03-scalability]{03} & Memory &
        Check memory stays practical as $N$ grows (engineering supplement). &
        --- \\
        \hyperref[subsec:study-04]{04} & CV parser &
        Measure CV parsing error before scoring. &
        1 \\
        \hyperref[subsec:study-05]{05} & JD parser &
        Measure JD parsing error before TF-IDF configuration. &
        1 \\
        \hyperref[subsec:study-06-adversarial]{06} & Adversarial &
        Stress-test fusion and integrity against stuffing, fraud, and gaming (RQ2 supplement). &
        --- \\
        \hyperref[subsec:study-07-spearman-high-discrimination]{07} & High disc. &
        Expert $\rho$: does MERIT Full lead when sources agree? &
        2 \\
        \hyperref[subsec:study-08-spearman-seniority]{08} & Seniority &
        Expert $\rho$: do engines over-rank seniority for an intern JD? &
        2 \\
        \hyperref[subsec:study-09-spearman-peer]{09} & Peer comp. &
        Expert $\rho$: does fusion break ties when CVs are similar? &
        2 \\
        \hyperref[subsec:study-10-spearman-dissonance]{10} & Dissonance &
        Expert $\rho$: does fusion fail when sources contradict the CV? &
        2 \\
        \hyperref[subsec:study-11-shapley-verification]{11} & Shapley &
        Check glass-box source attributions satisfy Shapley axioms. &
        3 \\
        \hyperref[subsec:study-12-bayesian-conflict-resolution]{12} & Bayes fusion &
        Check Beta--Bernoulli fusion favours demonstrated evidence over claims. &
        2 \\
        \hyperref[subsec:study-13-metric-prediction]{13} & TF-IDF &
        Check whether TF-IDF priorities match recruiter consensus. &
        1 \\
        \hyperref[subsec:hci-trial-protocol]{14} & HCI pilot &
        Check whether glass-box explanations change recruiter decisions and trust. &
        3 \\
        \hyperref[subsec:bias-mitigation-efficacy]{15} & Blind Mode &
        Check whether Blind Mode reduces proxy-driven shortlisting in the UI. &
        3 \\
        \hline
    \end{tabularx}
\end{table}

Table~\ref{tab:eval-literature-traceability} demonstrates how our literature review gaps tie directly into our evaluation studies and their success criteria. Within this mapping, temporal skill decay~\cite{arthur1998} is analysed in Section~\ref{subsec:study-08-skill-decay} as an RQ2 supporting analysis alongside the fusion studies. Additionally, Study~15 tests UI proxy debiasing~\cite{highhouse2008stubborn} rather than Shapley attributions, while Studies~02,~03, and~06 serve as engineering or adversarial supplements (Section~\ref{sec:fairness-computational-rigour}).

\begin{table}[H]
    \centering
    \caption[Literature-to-evaluation traceability]{Literature-to-evaluation traceability: mapping research questions to prior work, studies, and success criteria.}
    \label{tab:eval-literature-traceability}
    \footnotesize
    \renewcommand{\arraystretch}{1.45}
    \setlength{\tabcolsep}{6pt}
    \setlength{\extrarowheight}{2pt}
    \def\tabularxcolumn#1{m{#1}}%
    \begin{tabularx}{\textwidth}{@{}%
        >{\centering\arraybackslash}m{0.8cm}
        >{\centering\arraybackslash\hsize=1.05\hsize}X
        >{\centering\arraybackslash}m{1.9cm}
        >{\centering\arraybackslash\hsize=1.0\hsize}X
        >{\centering\arraybackslash\hsize=0.95\hsize}X @{}}
        \hline
        \textbf{RQ} & \textbf{Gap / prior work} & \textbf{Studies} & \textbf{Primary metrics} & \textbf{Success criterion} \\
        \hline
        RQ1 &
        Gugnani/Lo rigid weighting, ResumeBench parsing &
        04, 05, 13 &
        Skill $F_1$, name match, $\rho$ vs.\ recruiters &
        Parsing $\geq 80\%$, metric alignment exploratory \\
        RQ2 &
        Montandon supplementary GitHub, CV self-report bias &
        {\renewcommand{\arraystretch}{1}%
        \begin{tabular}[c]{@{}c@{}}12, 01, \\[-0.35ex] 07--10\end{tabular}} &
        Rank shift, Bayesian walkthrough, Spearman $\rho$ &
        $\rho \geq 0.7$ vs.\ experts (directional) \\
        RQ3 &
        Bogen/Raghavan transparency deficit, Lundberg Shapley &
        11, 14, 15 &
        Axiom checks, DRR, SUS, prestige gap &
        SUS $\geq 80$/100, blind gap compression \\
        \hline
    \end{tabularx}
\end{table}



\section{Evaluating Pillar 1: Extensibility}
\label{sec:foundational-ingestion-extensibility}

This section addresses \textbf{RQ1}: whether JD-driven, recruiter-configurable metrics 
and reliable ingestion can support downstream scoring.
In order for a screening system to be effective, it requires data to be ingested in a way that is accurate and reliable.
Therefore, the first area that we will evaluate is MERIT's ingestion layer, to ensure that it is reliable enough to
parse unstructured documents, specifically CVs and JDs. We do not consider structured data sources such as LinkedIn and GitHub
since there is no need to parse them, as they are already in a structured format. We begin with two studies, 
one for CVs and one for JDs, to evaluate the accuracy of the ingestion layer. Both studies compare the system's extraction 
outputs against human-verified ground-truth datasets, evaluating a total of 30 CVs (balanced across diverse layouts) and 15 JDs.
To finish the evaluation of this section, we then consider how well the metric priority predictions align with that
of a panel of recruiters, testing whether TF-IDF-derived priorities avoid the fixed-coefficient rigidity criticised in prior work 
by Gugnani and Misra~\cite{gugnani2020} and Lo et al.~\cite{lo2025aihiring}. Graphical representations for this section are 
in Appendix~\ref{appendix:ingestion-charts}.

\subsection{Evaluation Metrics}
\label{subsec:mathematical-metrics}

For our evaluation, we will be evaluating both set-based and string-based comparisons between the ground truth,
and the extracted data for our three core categories (skills, experience and identity).

\subsubsection{Skill Set Metrics}
Given $S_E$ as the set of extracted skills and $S_{GT}$ as the set of ground truth skills, we can
formulate the precision ($P_{skills}$), recall ($R_{skills}$), and F1-score ($F1_{skills}$) as follows:
{\small
\begin{align}
    P_{skills}  &= \frac{|S_E \cap S_{GT}|}{|S_E|}, &
    R_{skills}  &= \frac{|S_E \cap S_{GT}|}{|S_{GT}|}, &
    F1_{skills} &= \frac{2 P_{skills} R_{skills}}{P_{skills} + R_{skills}}
\end{align}
}

\subsubsection{Experience Recall}
For experience extraction, we mainly care about recall, since if we miss an experience, we might miss out on a
correct match. False positives are not a significant concern for this metric due to the strict nature of experience
extraction.

Given a set of ground truth experiences $E_{GT}$, and a set of extracted experiences $E_E$, we define a
successful match as follows:
\begin{equation}
    \text{Match}(e_{gt}) = \mathbb{I}\left(\exists e_e \in E_E \text{ s.t. } \text{norm}(e_{gt}) \subseteq \text{norm}(e_e)\right)
\end{equation}

Experience Recall ($R_{exp}$) is thus:
\begin{equation}
    R_{exp} = \frac{1}{|E_{GT}|} \sum_{e_{gt} \in E_{GT}} \text{Match}(e_{gt})
\end{equation}

\subsubsection{Identity and Metadata Matching}
When it comes to the extraction of identity for both CVs and JDs, we only consider the exact matches, since
metrics such as precision and recall are not applicable in that case. For CVs we do this for both name and 
email addresses. For the JD, while we do perform a similar operation, we are less strict with the matching,
it is considered a match if the ground truth label is a substring of the extracted label to account for the 
difficulty of matching the company name exactly in JDs.

\subsection{Study 04: CV Ingestion Accuracy}
\label{subsec:study-04}

We first evaluate the accuracy of the CV parser, extracting candidate data from 15 single-column CVs. We also
replicate this data across 15 multi-column CVs, ensuring that the content is the same for a given candidate,
to also evaluate the impact of a CV's layout on the accuracy of the parser. The CVs are designed to be diverse
in nature, across different roles and seniority levels. We craft a ground truth dataset for each candidate,
with the author verifying the accuracy and completeness of these JSON payloads.
Table~\ref{tab:cv_parser_results} presents the compiled results.

\begin{table}[htbp]
    \centering
    \caption[Study~04: CV ingestion by layout]{Study~04: CV ingestion accuracy by layout ($N=30$).}
    \label{tab:cv_parser_results}
    \merittablesize
    \setlength{\tabcolsep}{6pt}
    \begin{tabular}{l|c|c|c}
        \hline
        \textbf{Metric} & \textbf{Single-col.\ ($N{=}15$)} & \textbf{Multi-col.\ ($N{=}15$)} & \textbf{Overall ($N{=}30$)} \\
        \hline
        Skill Extraction Precision   & 69.6\%  & 69.6\%  & 69.6\%  \\
        Skill Extraction Recall      & 38.6\%  & 38.6\%  & 38.6\%  \\
        Skill Extraction $F_1$-Score & 49.4\%  & 49.4\%  & 49.4\%  \\
        Experience Recall            & 34.3\%  & 44.7\%  & 39.5\%  \\
        Name Match Rate              & 53.3\%  & 13.3\%  & 33.3\%  \\
        Email Match Rate             & 100.0\% & 100.0\% & 100.0\% \\
        \hline
    \end{tabular}
\end{table}

The distribution of individual scores across the template pool is plotted in
Figure~\ref{fig:ingestion-combined} (panel~a, Appendix~\ref{appendix:ingestion-charts}).

\subsubsection{Qualitative Analysis}

We have four key findings from this study, each with different root causes.

The first finding is that skill extraction is not hindered by the layout of a CV, but rather the dictionary that it is 
using to identify the skills. We can observe this as both layouts achieve the same precision, recall and F1-score ($69.6\%$, $38.6\%$ and $49.4\%$ respectively).
The source that we collect skills from, Devicon, catalogues programming languages as well as frameworks, however, it completely
ignores categories that are common on real engineering CVs. One such example is tooling for observability, such as Prometheus and Grafana.
The precision indicates that what is found is correct for the most part, but the limiting issue is the dictionary we use, rather than the 
logic of our extraction process. This finding is consistent with resume-parsing benchmarks such as ResumeBench \cite{emnlp2025resumebench}.

Another important observation is that the detection of candidate name drops sharply under multi-column layouts.
There is a 40\% drop between single-column and multi-column layouts when it comes to identifying the candidate name.
In a typical single-column layout, candidate name is generally the first item in the CV, at the top of the document.
However, for multi-column layouts, this is often not the case. This is a clear example of structural failure in the 
ingestion layer due to our parser having simplistic heuristics for name extraction.

We also find that email extraction is fully robust, achieving a 100\% match rate across both layout types,
confirming that our regex-based extraction process is not hindered by the layout of the candidate CV.

Finally, we end with counterintuitive results for experience extraction, noting that experience recall rises from
34.3\% in single-column layouts to 44.7\% in multi-column layouts. The structural difference between the two layouts
is definitely the reason, however, initially it was not clear why it was higher in the multi-column layout. Further
inspection leads us to believe that the reason is due to Single-Column layouts presenting experience as inline, with
multiple pieces of information on the same line, separated by a pipe character, or other delimiters. This must
be parsed as a single unit before we can isolate details such as the employer name and the role. Multi-column
layouts tend to have less width due to the introduction of more columns, so candidates tend to have these details on
separate lines. This actually makes it easier for the parser to isolate details, which could suggest why multi-column
layouts have better recall for experience compared to single-column layouts.

CVs are unstructured documents, and as such, there is a need for structured, schema-based data sources to be able
to account for parsing errors that may arise from CV ingestion (such as a LinkedIn profile). In real-world
applications, it is likely that commercial Applicant Tracking Systems (ATS) suffer from the same structural
parsing failures that we have observed in our study. The black-box nature of these screening systems would 
make it difficult for recruiters to know if this is actually the case for why a candidate receives a low score.
It is for this reason that candidates are strongly advised to use simple, single-column templates, as discussed in practitioner ATS guidance \cite{marr2019}.
Recruitment-management systems also tend to discard applicants who fail to mirror job-description keywords and criteria \cite{fuller2021hiddenworkers}.
In turn this uncovers a systemic bias: automated screening systems penalise highly qualified candidates if the
layout of their CV confuses the parser, consistent with the parsing failures we observed.

\textbf{Outcome:} Does not meet the internal parsing $\geq 80\%$ target (skill $F_1 = 49.4\%$, name match = $33.3\%$).
Precision is acceptable, but dictionary recall and multi-column name heuristics bound RQ1 ingestion.

\subsection{Study 05: Job Description Accuracy}
\label{subsec:study-05}
The JD parser was evaluated in a similar manner, with a total of 15 job descriptions across different software
engineering roles to ensure a wide range of formats and technical requirements. The compiled mean results
are summarised in Table~\ref{tab:jd_parser_results}, the per-document distribution is in
Figure~\ref{fig:ingestion-combined} (panel~b, Appendix~\ref{appendix:ingestion-charts}).

\begin{table}[htbp]
    \centering
    \caption[Study~05: JD extraction accuracy]{Study~05: mean JD extraction accuracy ($N=15$).}
    \label{tab:jd_parser_results}
    \merittablesize
    \setlength{\tabcolsep}{6pt}
    \begin{tabular}{l|l|r}
        \hline
        \textbf{Metric} & \textbf{Evaluation focus} & \textbf{Mean (\%)} \\
        \hline
        Skill Requirements Precision   & Target skill relevance       & 60.02\% \\
        Skill Requirements Recall      & Target skill coverage        & 34.13\% \\
        Skill Requirements $F_1$-Score & Overall requirement fidelity & 40.81\% \\
        Job Title Match Rate           & Role identification        & 46.67\% \\
        Company Match Rate             & Employer identification    & 6.67\%  \\
        \hline
    \end{tabular}
\end{table}

\subsubsection{Qualitative Analysis \& Strategic Ingestion Failure Cases}

Unlike CVs, which serve as concise professional summaries, Job Descriptions (JDs) are inherently
marketing-oriented documents designed to attract candidates. As a result, they are often dense with corporate
boilerplate, such as institutional missions, values, wellness benefits, and equal opportunity clauses.
This structural noise significantly degrades parsing accuracy. Because section headings are frequently
styled with generic tags (e.g.\ ``About Us'' or ``Who We Are''), the parser struggles to isolate core technical
qualifications from general corporate copy. This lexical contamination is directly reflected in the metrics,
driving a low Skill Recall of 34.13\% and restricting Skill Precision to 60.02\%, as the extraction engine
fails to separate critical technology requirements from marketing jargon.

The most notable observation within this evaluation is the poor performance of company name extraction, with a match rate
of 6.67\%, only successfully extracting one of the fifteen JD company names. The reason for this is structural, rather than 
semantics. Employers often represent company names in text that is not easily extractable by a parser, for example, they
represent their name as a logo or image. Text extractors (in our case, pdfminer) cannot extract this textual information from
an image, it is therefore forced to rely on the ``About Us'' section. If the name is not present in this section, or 
referred to via generic pronouns (such as ``Our company''). Then the extraction process fails. Fortunately, this is a non-critical
bottleneck, since the only usage of this metric is for the recruiter to know the company name, and not for the scoring process.
It is a convenience feature rather than a core functional requirement.

\textbf{Outcome:} Does not meet parsing $\geq 80\%$ (skill $F_1 = 40.81\%$).
While employer extraction is not critical for ranking, the JD skill recall still fails to meet RQ1's success criteria.

\subsection{Study 13: Metric Prediction vs Recruiter Alignment}
\label{subsec:study-13-metric-prediction}

While the past two studies evaluate MERIT's extraction, this study instead evaluates the weighting side of RQ1. We
aim to test whether TF-IDF derived metric priorities are able to align with expert consensus, as opposed to merely
having a fixed coefficient scheme.
We test \textbf{H1} (positive Spearman $\rho$ within each role) and \textbf{H2} (low MAE versus consensus on MERIT's 1--5 scale,
where 1 = essential and 5 = optional).

\subsubsection{Method}
When generating the market corpus, we employ Tonic.ai to synthesise a dataset of 100 JDs, providing a relatively realistic
dataset to aid in this evaluation. We have three hold-outs which are excluded from the generation of the corpus, that being a Junior 
Frontend, senior ML and Lead DevOps role. This is to prevent TF-IDF from being biased towards these roles. We have one global corpus,
we do not partition by seniority nor by role family.

On the testing side, we have two recruiters and two academic peers who are tasked with weighting the metrics given a JD. We intentionally
hide the suggestions that TF-IDF would have suggested to ensure fairness. Our consensus is determined by the median of these ratings, and 
we compare these results to TF-IDF suggestions, calculating the exact match, mean absolute error (MAE) and Spearman $\rho$.
The graphical representations for this study are available in Appendix~\ref{appendix:metric-prediction-charts}.

\subsubsection{Results}
As shown in Table~\ref{tab:study-13-metric-alignment}, TF-IDF is not well
calibrated for exact matches. We can see that none of the roles manage $\geq 50\%$ for the exact match rate between 
TF-IDF and recruiter consensus. However, we can observe that in most cases, Spearman $\rho$ is positive, showing some
alignment between the TF-IDF predictions and the recruiter consensus. This is a good sign, as it 
indicates that the TF-IDF predictions are generally in the right direction, but not necessarily a perfect match.
\begin{table}[htbp]
    \centering
    \caption[Study~13: TF-IDF vs consensus]{Study~13: TF-IDF vs.\ consensus ($N=4$, $n=27$ skill--role pairs).}
    \label{tab:study-13-metric-alignment}
    \merittablesize
    \setlength{\tabcolsep}{6pt}
    \begin{tabular}{l|c|c|c|c}
        \hline
        \textbf{Hold-out role} & \textbf{Skills} & \textbf{Exact match} & \textbf{MAE} & \textbf{Spearman $\rho$} \\
        \hline
        Junior Frontend (React) & 9  & 44.4\% & 0.78 & 0.52 \\
        Senior ML Engineer       & 8  & 25.0\% & 1.88 & $-0.45$ \\
        Lead DevOps / Platform   & 10 & 20.0\% & 1.20 & 0.40 \\
        \hline
        \textbf{Aggregate} ($n{=}27$) & 27 & 29.6\% & 1.26 & 0.25 \\
        \hline
    \end{tabular}
\end{table}

Of the 27 skill--role pairs, eight achieve exact agreement and a further nine differ by only one Likert band (63\% within $\pm 1$).
The near misses show that our metric predictions should be treated as reasonable guidance, rather than simply a default value.


\subsubsection{Limitations and their Implications}
\label{subsubsec:study-13-limitations}

This is a pilot study ($N = 4$ evaluators, $n = 27$ pairs), and the results are not generalisable to the wider market.
We are not able to conduct a larger panel study due to the time and resource constraints of this project.

Another limitation is due to the way we create the corpus, or rather, the fact that we only use one corpus for all TF-IDF calculations.
Due to this, we encounter cross-role IDF interference. To take an example, let's consider Python, a language used in many different roles such
as backend development, machine learning, data science and more. Since Python is common in these roles, it will inevitably be in a significant
proportion of job descriptions and, as a result, we would have a high IDF score, regardless of the role, making it seem like Python is 
always a common skill. This down-ranks Python on the metric predictions, regardless of the role.

In this study, we also did not ask our participants to spread their ratings evenly between the 1 and 5 scale. Compare this to TF-IDF
which has an in-built min--max normalisation to the range $[1.0, 5.0]$ which ensures a relatively even distribution of these ratings.
This disparity in rating distribution is likely to play a role in our evaluation results, particularly the exact match rate. 
That mismatch can depress exact-match rate and $\rho$ without invalidating the pilot, it is a
comparison caveat we report alongside the within-$\pm 1$ band counts above. Cross-study validity threats are summarised in
Section~\ref{sec:threats-to-validity}.

\textbf{Outcome:} Exploratory RQ1 criterion only: aggregate Spearman $\rho = 0.25$ with mixed hold-outs 
(Table~\ref{tab:study-13-metric-alignment}).
\textbf{H1} is partly supported and \textbf{H2} is not met. TF-IDF priorities should be treated as guidance, 
not automatic weights, with recruiter overrides on the frontend.

\subsection{Downstream Mitigation}
\label{subsec:downstream-mitigation}

These quantitative findings empirically justify two of MERIT's three core architectural pillars.

\subsubsection{Multi-Source Fusion (Pillar 2)}

The recall of 38.6\% means that CV-only screening will often miss the majority of a candidate's technical competencies
under the Devicon dictionary. The drawbacks of a dictionary-based approach are clearly visible here, as well as the 
reliance on a single source of truth for areas such as skill extraction. By ingesting other sources into candidate
scoring (such as GitHub repository language statistics and LinkedIn skill endorsements), MERIT can recover 
skills that the CV parser may have failed to recognise. 


\subsubsection{Glass-Box Explainability (Pillar 3)}
The issues presented above would often fail to be surfaced to recruiters in typical black-box scoring systems.
This highlights the need for a system that can surface these parsing failures to recruiters. The addition of 
MERIT's explainability layer achieves this by providing in-depth reviews of both the candidate's original 
data, as well as what MERIT was able to extract from it. This ensures that recruiters are able to audit and alter
their decisions based on the information that MERIT was able to gather. For example, a recruiter may observe that
a metric, such as Python, may not have been extracted from the candidate CV due to a parser error. The recruiter would
then be able to either observe the original CV to challenge the claim, or perhaps consider other data sources.
\section{Evaluating Pillar 2: Multi-Source Fusion}
\label{sec:evaluating-multi-source-fusion}

% Helper macros for evaluation tables (defined at top to prevent duplicate definition/undefined command errors)
\newcommand{\splithead}[2]{\begin{tabular}[c]{@{}c@{}}\textbf{#1} \\ \textbf{#2}\end{tabular}}
\newcommand{\splitcand}[2]{\begin{tabular}[c]{@{}l@{}}#1 \\ #2\end{tabular}}
\newcommand{\splitdesc}[2]{\begin{tabular}[c]{@{}l@{}}#1 \\ #2\end{tabular}}
\newcommand{\splitcell}[2]{\begin{tabular}[c]{@{}c@{}}\##1 \\ (#2)\end{tabular}}
\newcommand{\rankonly}[1]{\##1}

This section evaluates MERIT's multi-source fusion model, directly addressing RQ2. To ensure a thorough
evaluation, we proceed from the bottom up. We first demonstrate how skill decay can affect a candidate's score, 
from there, we then test the Bayesian Fusion process, evaluating the effects of conflict resolution and 
testing to see whether MERIT rewards demonstrated evidence over claimed evidence. From there, we conduct
an algorithmic study against baseline ATS scoring engines to compare alignment to expert decisions.

For studies involving ATS scoring engines, we specifically refer to four scoring engines. We have a \textbf{Traditional ATS}
baseline that replicated the keyword-matching system mentioned in Section~\ref{sec:automated-screening}. We also include a 
\textbf{Modern AI ATS} baseline that replicated the AI-assisted system mentioned in that section. For the evaluation
we also separate MERIT into two engines, one that only considers CV data, and one that considers all implemented
data sources (CV, GitHub, LinkedIn) to isolate the impact of the multi-source fusion layer.

The baselines are not commercial products, they are in-repository baselines that proxy keyword-matching and AI-assisted 
ATS systems. We intentionally omit rule-based engines as they are knockout-oriented, they do not rank candidates. 

We isolate the scoring logic by assuming perfect data extraction, using synthetic candidate JSON profiles. 
This removes parsing noise, establishing a clean mathematical baseline to accurately evaluate scoring logic, 
temporal decay, and conflict resolution. The validity implications for this are detailed
in Section~\ref{subsec:internal-validity}.

\subsection{Skill-Decay Sensitivity Analysis (RQ2 Supplement)}
\label{subsec:study-08-skill-decay}

Prior research on skill retention shows that unused workplace skills deteriorate over time~\cite{arthur1998}. 
The purpose of this study is to test whether our modelled skill decay provides a large enough mathematical 
separation between a candidate's active and stale skills to rank them appropriately. 

While this analysis does not strictly calibrate the decay constant ($\lambda$) against Arthur et al.'s empirical 
effect sizes, it models retention using the exponential decay curve defined in Equation~\ref{eq:exponential_decay}. 
Here, $S_0$ is the initial unweighted skill score and $\lambda$ is the decay constant. 

In MERIT, we set $\lambda \approx 0.25$ (yielding a half-life of $t_{1/2} \approx 2.77$ years, see 
Section~\ref{subsec:skill-decay-modelling}). To illustrate this 
sensitivity: for two otherwise identical profiles (active vs.\ stale), a 6-year stale signal yields 
$S(6) = S_0 e^{-0.25 \cdot 6} \approx 0.223 S_0$. Because the stale signal retains only 22.3\% of its active 
value, it creates a clear ranking separation. Consequently, the outcomes of this study should be seen as 
directional rather than strictly quantitative.

\subsection{Study 12: Bayesian Conflict Resolution}
\label{subsec:study-12-bayesian-conflict-resolution}

The reasoning for MERIT's multi-source fusion layer is to be able to reward candidates who are able to demonstrate their skills 
with publicly available evidence. This is particularly important for extensible metrics, where said metric is tied directly
to the job description that it was derived from. We believe it to be essential to reward candidates who can prove that they are
fit for the role, rather than just claim to be fit for it. In this study, we verify the Beta-Bernoulli conjugate updating 
implementation \cite{gelman2013bayesian} by contrasting a claims-only candidate against a demonstration-only candidate on the same extensible 
metric, motivated by the need to down-weight self-reported CV claims that personnel research shows are frequently 
embellished~\cite{henle2019_resume_fraud,dunning2004}. We also evaluate full-evidence and zero-evidence bounds to establish the 
mathematical ceiling and floor of the scoring engine.

In this walkthrough, the maximum signal that claims-only sources can achieve is $0.8$, while the maximum signal for demonstration-only signals
is uncapped at $1.0$. We start from priors $\alpha_0 = \beta_0 = 0.1$, and the step-by-step walkthrough for these two core profiles is given in Table~\ref{tab:bayesian-walkthrough}.

\begin{table}[htbp]
    \centering
    \caption[Study~12: Claims vs.\ demonstration walkthrough]{Study~12: Step-by-step fusion for the two core profiles.}
    \label{tab:bayesian-walkthrough}
    \small
    \setlength{\tabcolsep}{5pt}
    \begin{tabular}{l c c c c c c}
        \toprule
        \textbf{Step} ($i$) & $s_i$ & $c_i$ & $\alpha$ & $\beta$ & \textbf{Score} & $\sigma$ \\ \midrule
        \multicolumn{7}{c}{\textbf{Claimed only, no demonstration}} \\ \midrule
        0. Prior            & ---   & ---   & 0.100 & 0.100 & 0.500 & 0.456 \\
        1. CV               & 0.800 & 0.500 & 0.505 & 0.205 & 0.711 & 0.347 \\
        2. GitHub           & 0.000 & 0.900 & 0.506 & 1.106 & 0.314 & 0.287 \\
        3. LinkedIn         & 0.800 & 0.200 & 0.674 & 1.154 & 0.369 & 0.287 \\ \midrule
        \multicolumn{7}{c}{\textbf{Demonstrated only, no claims}} \\ \midrule
        0. Prior            & ---   & ---   & 0.100 & 0.100 & 0.500 & 0.456 \\
        1. CV               & 0.000 & 0.500 & 0.105 & 0.605 & 0.148 & 0.271 \\
        2. GitHub           & 1.000 & 0.900 & 1.006 & 0.606 & 0.624 & 0.300 \\
        3. LinkedIn         & 0.000 & 0.200 & 1.014 & 0.814 & 0.555 & 0.296 \\ \bottomrule
    \end{tabular}
\end{table}

Step-by-step parameter commentary for these core scenarios is in Appendix~\ref{appendix:bayesian-fusion-walkthrough}.
Table~\ref{tab:bayesian-scenario-comparison} subsequently summarises the final scores of these profiles alongside the empirical ceiling and floor bounds.
Figure~\ref{fig:bayesian-fusion-convergence} (Appendix~\ref{appendix:bayesian-fusion-charts}) plots the step-by-step trajectories of the two core scenarios.

\begin{table}[htbp]
    \centering
    \caption[Study~12: Claims vs.\ demonstration]{Study~12: Hypothetical claims vs.\ demonstration profiles, bounded by full and zero evidence.}
    \label{tab:bayesian-scenario-comparison}
    \small
    \setlength{\tabcolsep}{4pt}
    \begin{tabular}{l c c c c c}
        \toprule
        \textbf{Scenario} & $s_{\mathrm{CV}}$ & $s_{\mathrm{GH}}$ & $s_{\mathrm{LI}}$ &
        \textbf{Final} & $\sigma$ \\ \midrule
        Full evidence (Max claims and demonstration) & 0.8 & 1.0 & 0.8 & 0.861 & 0.206 \\
        Demonstrated only, no claims   & 0.0 & 1.0 & 0.0 & 0.555 & 0.296 \\
        Claimed only, no demonstration & 0.8 & 0.0 & 0.8 & 0.369 & 0.287 \\
        Zero evidence (No claims, no demonstration) & 0.0 & 0.0 & 0.0 & 0.062 & 0.144 \\ \bottomrule
    \end{tabular}
\end{table}

\textbf{The divergence point.} At the GitHub step (Table~\ref{tab:bayesian-walkthrough}, row~2, 
Figure~\ref{fig:bayesian-fusion-convergence}), the two core profiles diverge sharply. The 
demonstration-only candidate reaches $0.624$, while the claims-only candidate falls to $0.314$, a gap of 31 percentage points. This 
reflects the higher source confidence assigned to demonstrated GitHub evidence ($c_{\mathrm{GH}} = 0.9$), and the fact that a zero 
GitHub signal actively contradicts the CV claim rather than leaving it unchallenged.

\textbf{The claims-only penalty.} For the claims-only profile, a capped CV assertion ($s_1 = 0.8$, $c_1 = 0.5$) initially lifts the score from $0.500$ to $0.711$. This is
broadly the sort of outcome a CV-only screening pipeline might produce. If the candidate had simply not linked a GitHub profile (a missing source), 
the fusion engine would stop here and their score would remain at $0.711$, successfully avoiding unfair penalisation for lacking an open-source presence. 
However, because they linked a GitHub profile that contained no evidence for this specific metric, the zero signal actively contradicts the CV's claim. After GitHub contradicts the claim, the score crashes to $0.314$, and a
further LinkedIn endorsement ($s_3 = 0.8$, $c_3 = 0.2$) only recovers it to $0.369$. Additional self-report cannot substitute for missing
public work samples.

\textbf{The demonstration reward.} The demonstration-only profile follows the reverse pattern. An empty CV ($s_{\mathrm{CV}} = 0$) briefly pulls the score down to $0.148$
before GitHub ingestion, after which strong demonstrated evidence lifts it to $0.624$ and the final fused score settles at $0.555$. MERIT
rewards demonstrated evidence over claimed evidence, and this is an intentional design choice. Self-reported CVs are prone to embellishment,
while demonstrated evidence is harder to game, and requires a candidate to actually use the skill in their own projects or daily work.
The system favours candidates who can back up their skills with public work, while recording uncertainty when there is a lack of evidence.

\textbf{Deterring strategic omission.} One might ask whether a candidate could ``game'' the system by intentionally withholding a 
data source (like GitHub) to avoid the risk of a low signal contradicting their CV. However, because MERIT caps the strength and 
lowers the confidence of self-reported sources, a candidate who omits demonstrated evidence artificially ceilings their maximum 
possible score for that metric (and potentially all others). While their score will not actively crash, they leave the door wide 
open for other candidates with verified public work to easily surpass them in the final rankings. This inherently deters the strategic omission of data sources.

\subsection{Study 01: Comparative Engine and Source Ablation}
\label{subsec:study-01-comparative-ablation}

To evaluate the benefits and drawbacks of MERIT's multi-source fusion model, we compare it against three other baseline screening engines on
ten synthetic personas (\texttt{evaluation/01-comparative\_study/}).
The first is \textbf{Traditional ATS}, an in-repository engine that replicates the keyword-matching ATS category in
Section~\ref{sec:automated-screening}.
It scores candidates solely on the count of matching keywords in their CV.
The second is \textbf{Modern AI ATS}, an in-repository engine that replicates the AI-assisted ATS category.
It computes cosine similarities between CV embeddings and Job Description (JD) embeddings using Sentence-BERT (SBERT)~\cite{reimers2019sentence}.
To reiterate, neither baseline is a commercial ATS platform.
We also run \textbf{MERIT (CV Only)} which only considers CV data, and \textbf{MERIT (Full)} which considers all implemented data sources.

We can observe the rankings of the engines in Table~\ref{tab:algorithmic-ablation}, where rank~1 is considered the best
candidate--job fit and rank~10 the worst. The fusion shift column records how each candidate's rank changes from MERIT (CV Only) to MERIT (Full).
See Figure~\ref{fig:rank-displacement-bump-chart} (Appendix~\ref{appendix:rank_displacement}) for a bump-chart visualisation.


\begin{table}[H]
    \centering
    \caption[Study~01: Comparative benchmark rankings]{Study~01: ordinal ranks across four screening engines ($N=10$ personas, rank~1 = best fit).
    Fusion shift is the rank change from MERIT (CV Only) to MERIT (Full).}
    \label{tab:algorithmic-ablation}
    \vspace{1.5ex}
    \small
    \setlength{\tabcolsep}{5pt}
    \begin{tabular}{l l c c c c c}
        \toprule
        \textbf{Candidate} & \textbf{Profile} & \textbf{Trad.} & \textbf{Mod. AI} & \textbf{CV} & \textbf{Full} & \textbf{Fusion shift} \\ \midrule
        Alex Rivers        & Verified elite    & \rankonly{5}  & \rankonly{1}  & \rankonly{1}  & \rankonly{1}  & 0 \\
        Felix Vance        & Identity~mismatch & \rankonly{3}  & \rankonly{6}  & \rankonly{2}  & \rankonly{3}  & $-$1 \\
        Jordan Smith       & Hidden gem        & \rankonly{10} & \rankonly{9}  & \rankonly{10} & \rankonly{2}  & \textbf{+8} \\
        Buzz Ward          & Keyword~stuffer   & \rankonly{2}  & \rankonly{5}  & \rankonly{3}  & \rankonly{6}  & $-$3 \\
        Ghost Gary         & Unverified        & \rankonly{4}  & \rankonly{7}  & \rankonly{5}  & \rankonly{8}  & $-$3 \\
        Sam Old            & Stale profile     & \rankonly{6}  & \rankonly{8}  & \rankonly{7}  & \rankonly{5}  & +2 \\
        Vince Vault        & High-tenure Fraud & \rankonly{1}  & \rankonly{4}  & \rankonly{4}  & \rankonly{9}  & \textbf{$-$5} \\
        Fiona Frost        & Skill absence     & \rankonly{8}  & \rankonly{3}  & \rankonly{6}  & \rankonly{7}  & $-$1 \\
        Kim Junior         & Entry-level       & \rankonly{9}  & \rankonly{10} & \rankonly{9}  & \rankonly{4}  & +5 \\
        Carl Corp          & Closed-source     & \rankonly{7}  & \rankonly{2}  & \rankonly{8}  & \rankonly{10} & $-$2 \\ \bottomrule
    \end{tabular}
\end{table}

In the table, we can observe how the Traditional ATS prioritises keyword density and tenure (Vince Vault \#1 and Buzz Ward \#2),
while the Modern AI ATS prioritises CV semantics behind the keywords, rather than blindly prioritising density (Alex Rivers \#1 and Buzz Ward \#5).
Similar behaviour can be observed for CV-only MERIT, but note the lack of focus on tenure, specifically Carl Corp (\#2 $\rightarrow$ \#8 from Modern AI ATS to CV-only MERIT).
For MERIT (Full), we can observe that the multi-source data for Jordan Smith and Kim Junior are recognised, leading to rank shifts of +8 and +5 respectively.

\textbf{Jordan Smith: hidden gem.} The key finding from this study is from the hidden gem persona, Jordan Smith. The candidate is a low-tenure developer, only having
a year of formal experience on his CV, but has a strong GitHub presence that he has not mentioned well on his CV. While this is not a common scenario,
we mainly highlight this to show the cases where a candidate may not present themselves in an ATS-friendly manner, leading to them being overlooked by these engines.

MERIT (Full)'s integration of multi-source data allows the engine to take his extensive GitHub data into account, noting his public work in both Java and TypeScript,
two highly weighted skills in the JD. Due to the candidate having a significant amount of evidence compared to his peers, he is able to rise from rank \#10 to rank \#2.
This candidate is a clear example of how MERIT can reveal candidates who may be overlooked by current state-of-the-art single-source ATS engines that often
fall into the ``Resume Black Hole''~\cite{fuller2021hiddenworkers}.

\textbf{Carl Corp: proprietary data gap.} MERIT is a system that thrives on the presence of public data, candidates are encouraged to showcase their best work on their public profiles.
Carl Corp is a candidate who presents a trade-off that MERIT must face, that being candidates who are not able to publicly showcase their work.
Due to the candidate's commits being in private repositories, we are unable to determine his true skill set, leading to him being down-weighted by the engine,
since he is not able to benefit from the multi-source fusion layer to its full potential. Under MERIT (Full), he falls from rank \#2 to rank \#10.

Here, we present a ``Proprietary Data Gap'' scenario, showing that MERIT is more likely to favour a junior with public work, over an experienced candidate
with private data. MERIT trades off verifiable evidence for the accessibility of such evidence.

\textbf{Vince Vault: identity squatting.}\label{subsubsec:identity-mismatch-resolution} Due to the difficulty of gaming demonstrated evidence, MERIT prepares for the possibility of candidates who may
attempt to pass off another's work as their own. We refer to this act as ``identity squatting''.
In our study, we present Vince Vault as a candidate who performs this malicious action. Vince is a candidate who claims
ten years of experience and also an expert in Java, but it is likely that he is unable to show this evidence through his own public work.
Because of this, the candidate, in this scenario, submits a link to another candidate's GitHub account so that someone else's demonstrated work can support
his application. Traditional ATS and Modern AI ATS do not consider multi-source data, only relying on his CV claims, as such, 
he is able to rank highly on both engines by simply claiming to have the skills. The same observation can be made for MERIT (CV Only), 
where he is able to rank highly on CV text alone.

Once we account for multi-source data, MERIT (Full) detects a name difference between the candidate's CV and GitHub account, specifically, the GitHub
account belongs to another individual, ``Alice Smith''. The Integrity Engine penalises this mismatch, leading to Vince Vault falling from rank \#4 to rank \#9
between MERIT (CV Only) and MERIT (Full). To ensure fairness, we surface this information to the recruiter, providing the full names associated with the data sources
to the user to investigate further.

Identity squatting is rare in practice, but it is an important edge case to consider to ensure the integrity of the system.
\section{Evaluating Pillar 3: Glass-Box Explainability \& HCI}
\label{sec:evaluating-explainability}

This section evaluates the final pillar of MERIT, addressing RQ3. We first verify the mathematical correctness of the Shapley values,
applying the coalition Shapley formalism from explainable ML~\cite{lundberg2017unified} at the data-source level.
Then we conclude with a user study to understand the impact of the explainability layer on recruiter decision-making. In that study, we also
consider usability, confidence calibration and decision reversal rate for features such as Shapley values, audit trail and Blind Mode.
We evaluate mathematical correctness of Shapley attributions and the HCI impact of the Glass-Box layer on recruiters.

\subsection{Study 11: Shapley Value Verification}
\label{subsec:study-11-shapley-verification}

MERIT applies the Shapley formalism at the data-source coalition level 
(Section~\ref{sec:explainable-ai-shapley}), distinct from feature-level SHAP 
explanations for individual model predictions~\cite{lundberg2017unified}.
To ensure source-level attributions are mathematically sound, we verify Efficiency, Null Player, Additivity (per-metric), and Determinism 
on ten synthetic personas ($N = \{\text{CV}, \text{GitHub}, \text{LinkedIn}\}$). We intentionally omit the axiom of Symmetry since source
contributions for MERIT are hetrogenous by design, and therefore cannot produce identical marginal contributions in practice.

\subsubsection{Quantitative Verification Results}
Table~\ref{tab:shapley-verification} reports per-candidate results 
($\sum \phi_i = v(N)$ within floating-point tolerance, negative $\phi_i$ occur when a source triggers a penalty, 
e.g.\ GitHub identity mismatch).

\begin{table}[htbp]
    \centering
    \caption[Study~11: Shapley verification]{Study~11: per-candidate Shapley verification ($\sum \phi_i = v(N)$).}
    \label{tab:shapley-verification}
    \vspace{1.5ex}
    \small
    \setlength{\tabcolsep}{7pt}
    \begin{tabular}{l|ccc|cc}
        \hline
        \textbf{Candidate} & $\phi_{\mathrm{CV}}$ & $\phi_{\mathrm{GH}}$ & $\phi_{\mathrm{LI}}$ & $\sum \phi_i$ & $v(N)$ \\ \hline
        Alex Rivers   & +0.3107 & +0.4607 & +0.0087 & 0.7800 & 0.7800 \\
        Buzz Ward     & +0.3545 &  0.0000 & +0.0075 & 0.3620 & 0.3620 \\
        Carl Corp     & +0.1635 &  0.0000 & +0.0155 & 0.1790 & 0.1790 \\
        Felix Vance   & +0.3735 & +0.0660 & +0.0075 & 0.4470 & 0.4470 \\
        Fiona Frost   & +0.2660 & +0.0660 & +0.0080 & 0.3400 & 0.3400 \\
        Ghost Gary    & +0.2940 &  0.0000 &  0.0000 & 0.2940 & 0.2940 \\
        Jordan Smith  & +0.1028 & +0.5288 & +0.0113 & 0.6430 & 0.6430 \\
        Kim Junior    & +0.1340 & +0.2550 &  0.0000 & 0.3890 & 0.3890 \\
        Sam Old       & +0.1648 & +0.1878 & +0.0153 & 0.3680 & 0.3680 \\
        Vince Vault   & +0.2652 & $-$0.0708 & +0.0117 & 0.2060 & 0.2060 \\ \hline
    \end{tabular}
\end{table}

\begin{itemize}
    \item \textbf{Efficiency:} For all 10 candidates, $\sum \phi_i = v(N)$ within floating-point tolerance (Table~\ref{tab:shapley-verification}).

    \item \textbf{Null player:} Ghost Gary receives $\phi_{\text{GitHub}} = \phi_{\text{LinkedIn}} = 0$ on empty channels.

    \item \textbf{Determinism:} Five repeated runs produced zero drift (exact analytical Shapley, not Monte Carlo).

    \item \textbf{Per-metric efficiency:} Per-metric Shapley sums matched grand-coalition metric scores across all candidate--metric pairs.
\end{itemize}

For a visualisation of these results, refer to Figure~\ref{fig:shapley-combined} in 
Appendix~\ref{appendix:shapley-charts}.

\subsubsection{Key XAI Insights}

One key finding from this study is visible in the GitHub Shapley value for Vince Vault, where we show the presence of negative attributions.
In the context of cooperative game theory, this would occur when the introduction of a player leads to a reduction in the overall coalition's
winnings. For MERIT, this would occur when the introduction of a source leads to a fall in that candidate's match score. In this case,
it was due to the introduction of malicious intent via identity squatting, leading to a significant penalty being applied to the candidate's score.
Since Felix only provided false evidence for the GitHub source, negative attributions were focused there, rather than on LinkedIn or CV.
This penalty is then able to be presented to the recruiter to show that a candidate is allegedly claiming evidence that they do not have.
It is now the recruiter's responsibility to investigate the engine further to see if the fuzzy match failed due to actual intent of
malice, or due to a false positive.

Another finding can be presented through our hidden gem in Section~\ref{subsec:study-01-comparative-ablation},
Jordan Smith. As shown in that study, pivot from MERIT (CV Only) to MERIT (Full) boosts his score from rank 10 to rank 2. Due
to the introduction of Shapley values, we can truly confirm that this boost in rank is due to the inclusion of the GitHub data source,
which contributed $\phi_{GitHub} = 0.529$ to his score. This is presented to the recruiter on the frontend to show that the majority
of his score is demonstrated, not claimed. Once again, it is now the recruiter's decision to investigate the applicant further to
see if they are a good fit for the role, since they have demonstrated well that they align with the job, but their claims don't
seem to carry the same weight as their demonstrated evidence.

It is important to note that the dataset for this study reuses that of Study~01, therefore, we do not have significant Shapley values for
any candidate when it comes to the LinkedIn source attribution. This is due to the study being designed to evaluate the effect of
demonstrated evidence when it comes to candidate selection, so we kept LinkedIn evidence minimal to stress-test this CV-versus-GitHub contrast.
See Section~\ref{subsec:construct-validity} for more details. Note that this does not affect
the validity of this study in any way, we are primarily providing an explanation as to why LinkedIn doesn't contain significant
Shapley values.

\begin{table}[htbp]
    \centering
    \caption[Study~11: Aggregate Shapley by source]{Study~11: aggregate positive Shapley shares by source.}
    \label{tab:shapley-source-dominance}
    \vspace{1.5ex}
    \small
    \begin{tabular}{l|c|c}
        \hline
        \textbf{Source} & \textbf{Sum of positive $\phi_i$} & \textbf{Share of pool (\%)} \\ \hline
        CV (claimed)       & 2.429 & 59.6 \\
        GitHub (demonstrated) & 1.564 & 38.4 \\
        LinkedIn (claimed)    & 0.085 & 2.1 \\ \hline
    \end{tabular}
\end{table}

We can observe that in this cohort, CV evidence dominates the overall score, contributing to 59.6\% of the overall score of 
the candidates. This is followed by demonstrated GitHub evidence, contributing to 38.4\% of the overall score.
While LinkedIn evidence is present, it contributes to only 2.1\% of the overall score, indicating that it is not a significant
factor in the candidate's score. Note the margin for error since we only consider positive $\phi_i$ for the denominator (we exclude 
the negative GitHub attribution from Felix Vance).

\subsection{Study 14: HCI Trial}
\label{subsec:hci-trial-protocol}

Study 11 mathematically confirms that the explainable AI (XAI) engine is consistent, Study 14 evaluates the human-facing 
glass-box design against the transparency deficit identified in algorithmic hiring reviews~\cite{bogen2018,raghavan2020}.
Lo et al.~\cite{lo2025aihiring} report HR agreement on generated scores, but their published evaluation does not include a 
before/after test of whether recruiters reverse shortlist decisions once explanations are shown. 
We include the Decision Reversal Rate (DRR) metric to target this gap.
The purpose of our XAI is to provide transparent explanations to the recruiter, so we deem it important to evaluate
how our Glass-Box influences recruiter behaviour, hence our need for this metric.

In this study, we measure not only the decision, but also the confidence of recruiters with regards to their decision.
We also measure how usable this engine is, since if the system is too difficult to use, recruiters would likely avoid using it,
falling back into automation bias to avoid significant cognitive load. Therefore, we have the following hypotheses:

\begin{itemize}
    \item \textbf{Hypothesis 1 (Decision Calibration):} Access to the Glass-Box XAI audit trail will increase the rate at 
    which recruiters reverse incorrect screening decisions for adversarial candidates (such as
          keyword-stuffers and identity squatters) and "hidden gems" (like high-potential candidates with low tenure).
    \item \textbf{Hypothesis 2 (Confidence Calibration):} Recruiters will exhibit a higher alignment between their
          screening confidence and the empirical evidence, resulting in increased confidence for verified candidates and
          decreased confidence for unverified profiles.
    \item \textbf{Hypothesis 3 (Usability and Trust):} The interactive explainability layer will achieve a high System
          Usability Scale (SUS) score and receive high qualitative ratings for transparency and defensibility compared to a
          traditional black-box score.
\end{itemize}

\subsubsection{Two-Stage HCI Trial}
We start with a two-stage user study involving professional technical recruiters and hiring managers as participants.
This trial follows a two-stage process for each candidate in the pool:

\subsubsection{Stage 1: Control Condition (Black-Box Score)}
Recruiters are given a candidate CV with a single aggregate match alongside per-metric scores, but with no explanations.
Recruiters must make a decision on whether to shortlist or reject the candidate, and rate their confidence on a Likert scale
(1 = Not Confident at all, 5 = Extremely Confident). They must also provide a brief qualitative justification for their decision.

\subsubsection{Stage 2: Treatment Condition (Glass-Box Scorecard)}
Recruiters are now shown the same candidate, but with access to the glass-box scorecards, including the Shapley layer, audit trail,
and other explainability features. Recruiters will then be asked to state whether they wish to reverse or maintain their decision.
They are also asked to rate their confidence, along with an explanation of how the visualisations influenced their assessment, if at all.

\subsubsection{Key Quantitative Metrics and Empirical Results}
This pilot user trial includes 5 participants, consisting of 2 professional recruiters ($R_1, R_2$) with extensive experience in
software development talent acquisition, and 3 technical peers ($P_1, P_2, P_3$), MEng Computing students at Imperial College London, all of whom were acquainted with the author. 
These participants are familiar with the technical recruitment process and have undergone multiple technical projects and internships. 
Given $N=5$, we treat this as a formative pilot and report descriptive statistics, inferential tests are 
exploratory only (Section~\ref{sec:threats-to-validity}).

To prevent participant fatigue, we present each evaluator with only 4 synthetic candidates to evaluate. Unlike the personas
mentioned in previous sections, these candidates were designed in a real-world scenario. We used real candidate information
as the basis for synthesising these candidates:
\begin{itemize}
    \item \textbf{Moinecha Ramos-Horta} (Verifiable Elite Developer)
    \item \textbf{Yusuf Nilsson} (The "Hidden Gem" - 1 year of experience on CV, elite GitHub)
    \item \textbf{Robert Todd} (Adversarial Keyword Stuffer)
    \item \textbf{Aaron Mohammed} (Identity Squatter - claims deep Java/Go systems expertise on CV, squatted GitHub)
\end{itemize}

We present the quantitative results of this trial in Table~\ref{tab:hci-trial-results}.
Stage~1 and Stage~2 counts use $S$ (shortlist) and $R$ (reject), \textbf{DRR} is the fraction of Stage~1 decisions reversed after the Glass-Box scorecard.
Confidence uses a 5-point Likert scale (1 = not confident, 5 = extremely confident).
Decision-reversal, confidence-shift, and SUS distributions are in Figure~\ref{fig:hci-combined}
(panels a--c, Appendix~\ref{appendix:hci-trial-charts}). 
Raw session data are in \texttt{evaluation/14-hci\_trial/}.


\begin{table}[htbp]
    \centering
    \caption[Study~14: HCI pilot results]{Study~14: HCI pilot results ($N = 5$ evaluators).}
    \label{tab:hci-trial-results}
    \vspace{1.5ex}
    \small
    \setlength{\tabcolsep}{5pt}
    \begin{tabular}{l|c|c|c|c|c|c}
        \hline
        \textbf{Candidate}   & \textbf{Stage 1 (ctrl)} & \textbf{Stage 2 (trt)} & \textbf{DRR} & \textbf{$C_1$ (1--5)} & \textbf{$C_2$ (1--5)} & \textbf{$\Delta C$} \\ \hline
        Moinecha Ramos-Horta & $5S$ / $0R$            & $5S$ / $0R$            & 0.0\%        & 3.80                    & 4.60                    & +0.80                 \\
        Yusuf Nilsson        & $0S$ / $5R$            & $4S$ / $1R$            & 80.0\%       & 4.20                    & 4.80                    & +0.60                 \\
        Robert Todd          & $4S$ / $1R$            & $1S$ / $4R$            & 75.0\%       & 3.60                    & 4.40                    & +0.80                 \\
        Aaron Mohammed       & $3S$ / $2R$            & $0S$ / $5R$            & 100.0\%      & 3.20                    & 4.80                    & +1.60                 \\ \hline
    \end{tabular}
\end{table}

This pilot provides directional support for all three hypotheses under realistic hiring constraints, 
none are confirmed at inferential strength given $N=5$. We detail the results below:

\subsection*{Quantitative Results}

\subsubsection{Decision Calibration (Hypothesis 1)}
During the first stage, we can see the issue of the ``resume black hole'': all 5 evaluators rejected Yusuf Nilsson
due to his low tenure (1 year of experience on CV). However, upon transitioning to stage 2, we see that 4/5 evaluators
(80\% DRR) reversed their decision, recognising his valuable GitHub contributions. Recruiter $R_2$ chose to
maintain the rejection, noting that the JD strictly required 3+ years of experience regardless
of technical strength. This illustrates how human decision-makers calibrate system outputs with non-quantifiable
operational constraints.

Another key finding is that the explainability layer successfully aided in unmasking adversarial manipulation. More specifically,
we turn to the candidate \textbf{Robert Todd}, who managed to fool 4 out of 5 evaluators in stage 1 due to his keyword-dense
resume and aggregate score. However, upon transitioning to stage 2, we see that 3 out of 4 of those evaluators (75\% DRR)
reversed their decision, recognising his keyword stuffing.

The most significant instance of this is with the candidate \textbf{Aaron Mohammed}, who managed to fool 3 out of 5 evaluators
in stage 1 due to his gamified CV. However, upon transitioning to stage 2, rather than seeing supported evidence, recruiters were
presented with a significant red flag. Aaron Mohammed provided a GitHub profile that was not their own in order to pass off
someone else's work as his own. This is a serious ethical breach and a violation of professional integrity. Due to this, all
3 evaluators reversed their decision, rejecting Aaron Mohammed. It is crucial to note that recruiters did not reject
Mohammed based on the fact that he was not a good developer (in fact, according to his CV he was an expert in Java and Go).
Instead, they rejected him based on the fact that he showed a lack of integrity and dishonesty. This illustrates how the explainability layer
can be used to identify and reject participants who are not a good fit for the company culture, even if they are technically skilled.


\subsubsection{Confidence Calibration (Hypothesis 2)}
Once evaluators were presented with the Glass-Box MERIT interface, we observed a positive confidence shift ($\Delta C$)
across all profiles. This was particularly notable in the case of Moinecha Ramos-Horta, where her average confidence
rose from $3.80$ (Stage 1) to $4.60$ (Stage 2), an increase of $0.80$. In a similar fashion, confidence for Aaron Mohammed
also increased, rising from $3.20$ (Stage 1) to $4.80$ (Stage 2), an increase of $1.60$, which was
due to the veto penalty that stemmed from his dishonesty, giving the evaluators the confidence to reject him with high conviction.

The average confidence across all 20 profile-level observations (5 participants $\times$ 4 candidates) increased from $3.75$ 
(Stage 1) to $4.45$ (Stage 2), a mean increase of $0.70$ on a 5-point scale. A paired $t$-test on these rows yields 
$t(19) = -5.60$, $p < 0.001$, but the rows are not independent (fixed Stage~1$\rightarrow$Stage~2 order, repeated participants, 
acquaintance bias among peers). We therefore report this as descriptive evidence for Hypothesis~2, not confirmatory inference.

\subsubsection{Usability (Hypothesis 3)}
To evaluate the usability of MERIT's glass-box architecture, we followed the standard System Usability Scale (SUS) survey \cite{brooke1996sus}.
MERIT achieved an average SUS score of \textbf{81.50} out of 100, with individual scores ranging from 77.5 to 85.0.
According to the Bangor et al. adjective rating scale for SUS scores \cite{bangor2009_sus}, this corresponds to a ``Good'' rating (Grade B) at face value.
With $N=5$, it indicates interface feasibility only, not product-market validation (Section~\ref{subsec:construct-validity}).

\subsubsection{Qualitative themes}

The justifications presented by our evaluators can be found in \texttt{evaluation/14-hci\_trial/test\_data/responses.json}. As shown, there are
alignments with the quantitative results in Table~\ref{tab:hci-trial-results}. Take the evaluator
that still shortlisted Robert Todd after the keyword-stuffing alert. Despite what is considered a ``serious red flag'' by $P_3$, they still plan to 
``keep him shortlisted'' for a brief phone screen. Note how the implementation of the glass-box not only introduces contestability with the 
scores presented, but also offers recruiters the means to investigate claims made by candidates in better detail where evidence may be missing.

Meanwhile, other penalties such as identity squatting are treated as decisive when malicious dishonesty is surfaced. In the case of Aaron Mohammed,
recruiter $R_1$ states that this is an ``immediate reject'' due to the ``Identity Mismatch Veto'' penalty upon seeing that the GitHub account's
full name does not match the candidate's name. Passing off another's work as one's own is a serious breach of honesty and integrity, and is often 
not a good fit for company culture. 

The responses point towards cases in which automation bias can be challenged by a recruiter's human judgement. While this study confirms the 
theoretical success of the Glass-Box design, we must note that this is a pilot study. Studies with a larger recruiter panel would be needed
before stronger claims can be made.

\textbf{Outcome:} Meets the internal SUS target (81.5/100).
Decision-reversal and confidence shifts provide directional support for glass-box calibration, but the pilot is not confirmatory ($N=5$).
\section{Human-Machine Alignment}
\label{sec:human-machine-alignment}

Study~01 in section~\ref{subsec:study-01-comparative-ablation} 
included simplistic adversarial personas to stress-test rank shifts on specific candidate characteristics to see
how each engine treats their specific nuances when it comes to candidate selection. Candidates were not designed to
be representative of a real-world candidate. Here we use synthetic cohorts built from realistic candidate profiles and compare
engine rankings against recruiter consensus, using the same baselines as Study~01.

We compare these engine rankings across four synthetic cohorts of size 10, using the Spearman rank correlation coefficient ($\rho$).
This follows Lo et al.'s~\cite{lo2025aihiring} comparison of automated screening output against human reference rankings. We extend
this with ATS baselines and multi-source scoring. Spearman $\rho$ measures agreement between two orderings only, not hiring quality.
For graphical representations of the results, refer to Figure~\ref{fig:spearman-studies-combined} 
in Appendix~\ref{appendix:spearman-charts}.


\subsection{Methodology}
\label{subsec:spearman-method}

Each study pairs a set of ten synthetic candidates with a job description.
For automated scoring, the cohort includes GitHub and LinkedIn signals where the study design requires them, 
and supplementary evidence agrees with each candidate's profile except in Study~10, where we purposely create 
contradictory GitHub evidence to test fusion robustness.
Recruiter consensus, however, is formed from CV PDFs only (see below).
Similar to the earlier studies, rank~1 is the best fit.

\textbf{Recruiter consensus protocol.}
The reference order for each cohort was produced by the same four professional technical recruiters ($R_1$--$R_4$) across Studies~07--10.
Each recruiter independently reviewed the role JD and an anonymised bundle of ten candidate CV PDFs.
They did not receive GitHub or LinkedIn profiles, only the material a recruiter would typically see at CV-screening stage.

We aggregated the four lists as follows.
First, we computed each candidate's mean rank across recruiters (lower is better).
When two or more candidates tied on mean rank, we used the lower median rank.
Any remaining ties were resolved by panel discussion until all four recruiters agreed on the final order.
The resulting order is stored in \texttt{human\_rankings.txt} and summarised in \texttt{reasoning.md} under each 
study directory in \texttt{evaluation/}.
We refer to this as expert consensus in the results below.
Spearman $\rho$ measures how closely each engine matches that CV-only recruiter consensus.
Engines may use additional GitHub and LinkedIn signals (MERIT Full), so agreement is not expected to be perfect in every cohort.

For each engine we produce a score-based ordering and compute $\rho$ against the recruiter consensus ranking.
\begin{equation}
    \rho = 1 - \frac{6 \sum_i (h_i - e_i)^2}{n(n^2 - 1)} .
\end{equation}

We use \texttt{scipy.stats.spearmanr} to compute the $\rho$ and two-sided $p$-values. Note that due to the small sample size
($n=10$), a single rank swap can move $\rho$ sharply, for a similar reason, the $p$-values are exploratory. To reiterate,
the engines we include are the \textbf{Traditional ATS} (keyword CV baseline), \textbf{Modern AI ATS} (a minimal in-repository
semantic baseline: SBERT skill matching plus experience weighting on CV text only), \textbf{MERIT (CV Only)}, and \textbf{MERIT (Full)}
(CV plus GitHub and LinkedIn). Neither ATS baseline is a commercial product. The Modern AI baseline models one stage of AI-assisted
screening and serves as a minimal proxy. We acknowledge that real-world commercial platforms (e.g., Eightfold or Workday) likely employ
much more sophisticated proprietary parsing and semantic matching than our baseline suggests, which may understate their true capabilities 
when comparing against MERIT (Section~\ref{subsec:internal-validity}).

\subsection{Spearman Studies 07--10}
\label{subsec:study-07-spearman-high-discrimination}

We present a summary of our findings in Table~\ref{tab:spearman-studies-combined}.
Study~07 includes a high-discrimination cohort for a full-stack backend engineer role, where candidates are diverse in 
terms of previous experience and seniority. Study~08 is a seniority-bias test, where all candidates are qualified for an 
intern backend engineer role, but they are of differing seniority. Study~09 is a peer competition test, where
candidates are of the same seniority but different engineering specialties. Study~10 is a signal dissonance case that replicates Study~07 with
deliberately contradictory GitHub evidence to test the robustness of the fusion engine.

\begin{table}[htbp]
    \centering
    \caption[Spearman $\rho$ across Studies 07--10]{Spearman $\rho$ by engine and cohort ($n=10$ per study). $p$-values are exploratory, bold values denote the highest $\rho$ in each cohort column.}
    \label{tab:spearman-studies-combined}
    \small
    \setlength{\tabcolsep}{4pt}
    \begin{tabular}{l|c|c|c|c}
        \hline
        \textbf{Engine} &
        \textbf{Study 07} &
        \textbf{Study 08} &
        \textbf{Study 09} &
        \textbf{Study 10} \\
        & High disc. & Seniority & Peer & Dissonance \\
        \hline
        Traditional ATS & 0.5515 ($p{=}0.098$) & $-$0.6606 ($p{=}0.038$) & \textbf{0.7939 ($p{=}0.006$)} & \textbf{0.5515 ($p{=}0.098$)} \\
        Modern AI ATS   & 0.2848 ($p{=}0.425$) & $-$0.8061 ($p{=}0.005$) & 0.2364 ($p{=}0.511$) & 0.2848 ($p{=}0.425$) \\
        MERIT (CV Only) & 0.5152 ($p{=}0.128$) & \textbf{$-$0.5152 ($p{=}0.128$)} & 0.0909 ($p{=}0.803$) & 0.5152 ($p{=}0.128$) \\
        MERIT (Full)    & \textbf{0.7333} ($p{=}0.016$) & \textbf{$-$0.5152 ($p{=}0.128$)} & 0.7212 ($p{=}0.019$) & 0.0909 ($p{=}0.803$) \\
        \hline
    \end{tabular}
\end{table}

\textbf{Note:} Due to the micro-cohort size ($n=10$), $p$-values are highly volatile and provided for transparency only,  
they do not imply statistical significance. We should not draw any conclusions about the generalisability of the results based on the $p$-values.

\textbf{Study 07 (high discrimination).}
MERIT (Full) achieves the strongest expert alignment ($\rho=0.733$, $p=0.016$), outperforming Traditional ATS ($\rho=0.552$),
MERIT (CV Only) ($\rho=0.515$), and the semantic baseline ($\rho=0.285$).
Recruiters ranked from CV PDFs alone, while MERIT (Full) also ingests GitHub and LinkedIn, raising $\rho$ from 0.515 to 0.733 
when that supplementary evidence agrees with the CV.

\textbf{Outcome:} Meets the directional $\rho \geq 0.7$ criterion for MERIT (Full) in this cohort ($\rho = 0.733$, $p = 0.016$).

\textbf{Study 08 (seniority bias).}
\label{subsec:study-08-spearman-seniority}
All engines fail to align with the expert consensus, achieving a negative $\rho$, ranking the principal/staff CVs above 
the intuitive junior candidate. Here, the impact of multi-source fusion is minimal, since multi-source fusion adds 
magnitude to each candidate relatively equally, leading to no rank changes. This is a stress test to show that 
multi-source fusion does not consider nuanced scenarios such as overqualified candidates.
\footnote{Note that the $\rho$ value of $-0.5152$ is exactly the inverse of the $0.5152$ observed in Studies~07 and 10. 
This is not a clerical error, but a pure mathematical coincidence. Due to the small cohort size of $n=10$, the Spearman 
correlation can only take on highly quantised fractions (in this case, $\pm17/33 \approx \pm0.51515$). The sum of squared 
rank differences happened to be exactly 80 in Study~07, and 250 in Study~08.}

\textbf{Outcome:} Does not meet $\rho \geq 0.7$ for any engine (MERIT Full $\rho = -0.515$); seniority-bias stress test failed.

\textbf{Study 09 (peer competition).}
\label{subsec:study-09-spearman-peer}
Traditional ATS aligns best with $\rho=0.794$, followed by MERIT (Full) with $\rho=0.721$ ($p=0.019$). Note the difference in $\rho$ between
MERIT (Full) and MERIT (CV Only), which is due to GitHub and LinkedIn evidence available to the engine but not shown to recruiters.
Multi-source fusion can act as a tie-breaker when CV PDFs alone are insufficient.

\textbf{Outcome:} MERIT (Full) meets $\rho \geq 0.7$ ($\rho = 0.721$) but does not lead the column (Traditional ATS $\rho = 0.794$); directional support for fusion as a tie-breaker only.

\textbf{Study 10 (signal dissonance).}
\label{subsec:study-10-spearman-dissonance}
The failure of MERIT (Full) is presented through the replication of Study~07, only with the intentional sabotage of
GitHub candidate evidence to contradict the CV narrative presented for a given candidate.
Recruiter consensus is unchanged from Study~07 because the CV PDFs are identical and recruiters never saw the sabotaged GitHub payloads.
This deliberate failure was 
orchestrated to maximise benefits of multi-source fusion for the worst performing candidates while minimising benefits
for the best performing candidates. This scenario demonstrates that multi-source fusion can be vulnerable to data poisoning,
however, note that it is uncommon for a candidate to intentionally contradict their CV with public repository evidence.
This case was selected to show that fusion performance is dependent on the quality of the evidence that it is supplied with,
it will underperform when sources present contradictory evidence.

\textbf{Outcome:} Does not meet $\rho \geq 0.7$ for MERIT (Full) ($\rho \approx 0.09$); deliberate adversarial failure under contradictory supplementary evidence.

\subsection{Synthesis}
\label{subsec:spearman-synthesis}
No engine dominates every scenario. MERIT (Full) is strongest when cohorts are diverse with consistent evidence (Study~07),
useful as a tie-breaker among peers (Study~09), but fails under batch-wide dissonance (Study~10). CV-only and full MERIT share
seniority-bias vulnerability with the baselines (Study~08). The keyword baseline can match experts more closely than our minimal
semantic baseline in some cohorts, however, MERIT's advantage is not ``more AI on the CV,'' but evidence-aware, job-weighted fusion when
sources align.

\section{Fairness, Security \& Computational Rigour}
\label{sec:fairness-computational-rigour}

The previous sections covered the accuracy of MERIT, its fusion logic, explainability and also its alignment with human experts.
We now move on to the last stage of our evaluation by testing the fairness, security against adversarial inputs and computational 
scalability of MERIT. Notably, the bias mitigation evaluation (Study~15) presented here serves as the final component of our RQ3 analysis, addressing the debiasing effects of the UI. Refer to Appendix~\ref{appendix:fairness-scalability-charts} 
and Appendix~\ref{appendix:bias-anonymisation-charts} 
for the charts associated with this section.

\subsection{Study 15: Bias Mitigation}
\label{subsec:bias-mitigation-efficacy}


Manual screening is prone to bias and heuristics as mentioned in Section~\ref{sec:background-manual-screening}.
Study~15 evaluates whether our Blind Mode interface alters recruiter decisions, specifically reducing the reliance on proxy markers such
as (prestige, layout and titles). This is a UI debiasing pilot study, it is not a fairness certificate.
Throughout the study, evaluators did not see the ranks for the candidates, but they did see criteria such as match score and 
Glass-Box breakdown. We follow the between-subjects design, making sure that evaluators only see one UI condition, making sure that
recruiter perception is not affected by a previous visible review of the same people.

This pilot study is guided by the following hypotheses:

\textbf{H1 (Proxy-gap compression):} The shortlist-rate gap between prestige-high and prestige-mid/low proxy groups will be smaller in the Blind arm than in the Visible arm.

\textbf{H2 (Title and layout debiasing):} Candidates penalised by intern titles, noisy parsing, or role-skewed visible layout (e.g. Vincent) will have higher shortlist rates in the Blind arm, with MERIT scores held constant.


\subsubsection{Experimental protocol}
For this study, we have a total of eight participants (four recruiters and four peer evaluators denoted as $RN$ and $PN$ respectively, where $N \in \{1, 2, 3, 4\}$).
We have two phases to the study, so to ensure fairness we split participants into two blocks that each contain two recruiters and two peers.

\textbf{Block 1: Visible arm} ($n = 4$, B1 $\in$ \{\textbf{R1, R2, P1, P2}\}): Intelligence Report popup with Blind Mode \textbf{off} (names, employers, institutions, titles, standard layout).

\textbf{Block 2: Blind arm} ($n = 4$, B2 $\in$ \{\textbf{R3, R4, P3, P4}\}): same popup with Blind Mode enabled (candidate identifiers, redacted institutions with tier retained, neutralised skill layout. Section~\ref{sec:bias-mitigation-safeguards}).

In this study, all participants review the same four profiles in a fixed order against the Intern Backend Engineer JD from Study~07.
We always hide the ranking page from evaluators, only showing the report popup. Due to our intentional design choices, match scores,
audit trails and confidence scores are always visible in the popup. After viewing the popup, evaluators record their decision (Shortlist/Reject),
and also their confidence in that decision (1--5) alongside a brief justification. We do not present evidence beyond the CV. No
evaluator saw both Visible and Blind Mode on the same candidate to prevent memory effects.
We have that $n = 4$ per arm due to the limited number of participants, which increases the variance of the results
but allows us to draw relatively reliable conclusions. Results are directional only due to the small sample size.

\subsubsection{Cohort design}
Synthetic candidates were selected from Studies~07 and~08, with no overlap with the Study~14 set 
(Moinecha Ramos-Horta, Yusuf Nilsson, Robert Todd, Aaron Mohammed):
\begin{itemize}
    \item \textbf{Aisling Hunt} (Study~07) --- junior data scientist, UC Berkeley, feminine name and gendered extracurricular signals on the visible CV.
    \item \textbf{Vincent Germain} (Study~08) --- intern backend engineer, human rank~1 in Study~08 but intern title and corrupted PDF header on the visible view.
    \item \textbf{Riko Al-Tayer} (Study~08) --- senior engineer, Cambridge/Imperial/Stripe pedigree (prestige-halo risk, overqualified for an intern JD).
    \item \textbf{Jude Washington} (Study~07) --- junior iOS developer, Berkeley education but unknown employers (TechFlow, Momentum), and a mobile-first visible layout.
\end{itemize}

As expected, MERIT Full scores are unchanged by Blind Mode (UI-only redaction): Riko (60.9\%) $>$ Jude (39.9\%) $>$ Aisling (36.7\%) $>$ Vincent (31.1\%).

\subsubsection{Quantitative results}
Table~\ref{tab:bias-audit-summary} reports how often each candidate was shortlisted in each arm ($n = 4$ evaluators per arm). 
Figure~\ref{fig:bias-combined} 
(panels a--b, Appendix~\ref{appendix:bias-anonymisation-charts}) 
plots the same data. Table~\ref{tab:bias-disparity-metrics} 
summarises the prestige proxy gap used for H1.

\begin{table}[htbp]
    \centering
    \caption[Study~15: Blind Mode shortlist rates]{Study~15: shortlist rates by candidate and UI arm ($n = 4$ per arm).}
    \label{tab:bias-audit-summary}
    \vspace{1.5ex}
    \small
    \begin{tabular}{l c c c c}
        \toprule
        \textbf{Candidate} & \textbf{MERIT (\%)} & \textbf{Visible (\%)} & \textbf{Blind (\%)} & \textbf{$\Delta$ (pp)} \\ \midrule
        Riko Al-Tayer      & 60.9 & 100.0 & 75.0 & $-$25.0 \\
        Aisling Hunt       & 36.7 & 25.0  & 75.0 & +50.0 \\
        Vincent Germain    & 31.1 & 0.0   & 75.0 & +75.0 \\
        Jude Washington    & 39.9 & 25.0  & 50.0 & +25.0 \\ \bottomrule
    \end{tabular}
\end{table}

\begin{table}[htbp]
    \centering
    \caption[Study~15: Aggregate disparity metrics]{Study~15: prestige proxy gap (high: Riko, Aisling; mid/low: Vincent, Jude).}
    \label{tab:bias-disparity-metrics}
    \vspace{1.5ex}
    \small
    \begin{tabular}{l c c}
        \toprule
        \textbf{Metric} & \textbf{Visible arm} & \textbf{Blind arm} \\ \midrule
        Evaluators (recruiters / peers) & 2 / 2 & 2 / 2 \\
        Mean shortlist rate --- prestige-high group (\%) & 62.5 & 75.0 \\
        Mean shortlist rate --- prestige-mid/low group (\%) & 12.5 & 62.5 \\
        Prestige gap (high minus mid/low, pp) & 50.0 & 12.5 \\ \bottomrule
    \end{tabular}
\end{table}

\paragraph{Visible arm}
All of the evaluators shortlist Riko, and all recruiters reject Vincent, noting his short tenure and corrupted header.
Aisling and Jude each receive one shortlist decision. In the visible arm, it is strongly suggested that decisions are
led by visible cues such as university names, employer logos and job titles, as opposed to the candidate-job fit.

\paragraph{Blind arm}
The prestige gap falls from 50pp to 12.5pp, as shown in Table~\ref{tab:bias-disparity-metrics}.
The two candidates leading this change are Vincent (0\% $\rightarrow$ 75\%) and Aisling (25\% $\rightarrow$ 75\%), as they were shortlisted
more often due to the hidden prestige cues. This supports H1 (Proxy-gap compression). H2 is also addressed through these two candidates,
since they were also often penalised for visible layout and titles. Another interesting observation is that, similar to Section~\ref{subsec:study-08-spearman-seniority},
overqualification is still a factor to consider since Riko was rejected by a recruiter for this very reason. Blind Mode does not
make everyone agree with MERIT's ordering, and it shouldn't be expected to. Blind Mode is designed to change who gets a fair first look
once we hide the names of employers and academic institutions.

\subsubsection{Study 15 takeaway}
Blind Mode is useful as a first-pass safeguard, it reduces obvious proxy bias in this pilot and helps profiles like 
Vincent get considered when titles and formatting had put them at a disadvantage. It does not replace Study~14's Glass-Box 
checks, and it is not proof of demographic fairness.

\textbf{Outcome:} Directionally meets the blind gap-compression criterion (prestige gap 50~pp $\rightarrow$ 12.5~pp).
Formative pilot only ($n = 4$ per arm), not inferential proof of fairness. It supports the direction of results, but that is not a
guarantee of fairness, see Section~\ref{sec:threats-to-validity} for further clarification.

\subsection{Study 06: Adversarial Stress Testing}
\label{subsec:study-06-adversarial}

Practitioner guidance documents keyword optimisation and layout gaming against ATS filters~\cite{marr2019,fuller2021hiddenworkers}, 
and personnel research documents intentional resume misrepresentation~\cite{henle2019_resume_fraud}.
Study~06 does not replicate those field studies, instead it stress-tests whether MERIT's integrity and fusion layers respond when 
synthetic personas emulate those attack patterns (keyword stuffing, metric inflation, identity squatting).

This study provides a small-scale evaluation of MERIT's reaction to a base candidate that has been duplicated with 
different perturbations to simulate different real-world scenarios and adversarial attacks. All candidates in this test are 
a copy of the same base candidate, \textbf{Honest Candidate}, and the perturbed candidate's name is often derived from
the scenario or attack they represent. \textbf{Honest Candidate} is a junior software engineer with a strong CV and a fair 
match percentage of 49.8\%. We have the following scenarios:
\begin{itemize}
    \item \textbf{Ghost}: Keyword stuffs their CV, but provides a public repository that contradicts their claims.
    \item \textbf{Fraud}: Inflates their CV claims.
    \item \textbf{Stale}: Provides outdated public-repository evidence.
    \item \textbf{Gamer}: Manipulates their CV and public-repository to artificially boost specific metrics.
    \item \textbf{Squatter}: Squats the name of another candidate. Caught by MERIT's Integrity Layer.
    \item \textbf{Smart Squatter}: Squats the name of another candidate. Fuzzy name-match failed to detect the attack.
    \item \textbf{Shadow}: Provides no public-repository evidence to support their claims (as opposed to Ghost, who has contradictory evidence).
    \item \textbf{Inflater}: Inflates their repository metrics only for one language with high LOC, and completely disregards others.
\end{itemize}

Each scenario is scored by both the CV-only and Full MERIT engines. 
Results are exported to \texttt{output/multi\_source\_adversarial\_matrix.csv}. 
Table~\ref{tab:adversarial-matrix} summarises the scores.
$\Delta_{\mathrm{honest\_cv}}$ is CV-only minus the honest CV-only score, $\Delta_{\mathrm{fusion}}$ is Full minus CV-only on the same row,
$\Delta_{\mathrm{honest\_fusion}}$ is Full minus the honest Full score 
(panel~b of Figure~\ref{fig:adversarial-combined}).
Figure~\ref{fig:adversarial-combined} 
(panels a--b, Appendix~\ref{appendix:fairness-scalability-charts}) 
compares CV-only (grey) versus
Full MERIT (blue) per scenario and reports deviation from the honest Full baseline (49.8\%, dashed line).

\begin{table}[htbp]
    \centering
    \caption[Study~06: Adversarial robustness matrix]{Study~06: adversarial robustness vs.\ honest baseline (49.8\% Full).}
    \label{tab:adversarial-matrix}
    \vspace{1.5ex}
    \footnotesize
    \resizebox{\textwidth}{!}{%
        \begin{tabular}{l l >{\centering\arraybackslash}c >{\centering\arraybackslash}c >{\centering\arraybackslash}c >{\centering\arraybackslash}c >{\centering\arraybackslash}c}
            \toprule
            \textbf{Scenario} & \textbf{Persona} & \textbf{CV (\%)} & \textbf{$\Delta_{\mathrm{honest\_cv}}$} & \textbf{Full (\%)} & \textbf{$\Delta_{\mathrm{fusion}}$} & \textbf{$\Delta_{\mathrm{honest\_fusion}}$} \\ \midrule
            Honest            & Honest Candidate   & 29.6 & $+0.0$ & 49.8 & $+20.2$ & $+0.0$ \\
            Ghost             & The Ghost          & 26.7 & $-2.9$ & 34.8 & $+8.1$  & $-15.0$ \\
            Fraud             & The Fraud          & 44.3 & $+14.7$ & 47.4 & $+3.1$  & $-2.4$ \\
            Stale             & Stale Candidate    & 35.0 & $+5.4$  & 37.9 & $+2.9$  & $-11.9$ \\
            Gamer             & The Time Traveler  & 35.0 & $+5.4$  & 58.1 & $+23.1$ & $+8.3$ \\
            Squatter          & Vince Vault        & 29.6 & $+0.0$  & 29.8 & $+0.2$  & $-20.0$ \\
            Smart Squatter    & Alex Rivers Smith  & 29.6 & $+0.0$  & 49.8 & $+20.2$ & $+0.0$ \\
            Shadow            & The Shadow         & 29.6 & $+0.0$  & 30.4 & $+0.8$  & $-19.4$ \\
            Inflater          & The Inflater       & 29.6 & $+0.0$  & 46.1 & $+16.5$ & $-3.7$  \\ \bottomrule
        \end{tabular}%
    }
\end{table}





\subsubsection{Successful Behaviours}
\textbf{Squatter.} This key highlight shows the Integrity Engine successfully detecting and neutralising identity squatting. Fusion adds an insignificant $+0.2$ uplift. The $\Delta_{\mathrm{honest\_fusion}}$ reveals a $20$-point deficit compared to an identical honest candidate, proving the Integrity Engine effectively cuts off the attacker's advantage.

\textbf{Stale.} This candidate updates their CV but provides outdated repository evidence. While CV-only is unchanged, fusion adds only a minor $+2.9$ uplift. Relative to the honest candidate, they fall $-11.9$ points behind. They are not explicitly penalised, rather, Bayesian Fusion naturally reflects the lower signal strength of their time-decayed GitHub evidence.

\textbf{Ghost.} This candidate keyword-stuffs their CV while providing a public-repository that contradicts their claims. CV-only scoring partially penalises this, but their $\Delta_{\mathrm{honest\_fusion}}$ score drops heavily to $-15.0$ because Bayesian Fusion updates the Beta Distribution to reflect the contradictory GitHub signal.

\textbf{Shadow.} Unlike Ghost, this candidate does not actively contradict their CV, but their complete lack of GitHub evidence prevents them from gaining the Bayesian boost awarded to the honest candidate. While they receive a minor $+0.8$ uplift from LinkedIn, their $\Delta_{\mathrm{honest\_fusion}}$ of $-19.4$ demonstrates that missing evidence safely leaves candidates behind rather than actively penalising them.

\textbf{Inflater.} This candidate inflates their metrics by concentrating massive lines of code (LOC) into a single language, disregarding others. Because MERIT evaluates all JD requirements, neglected languages receive no demonstrated evidence. This dilutes their overall probability of competence during the Bayesian update, dragging their $\Delta_{\mathrm{honest\_fusion}}$ down to $-3.7$.

\subsubsection{Partial Success}
\textbf{Fraud.} This candidate successfully inflates their CV-only score ($+14.7$) through major false claims, while providing contradictory repository evidence. Unlike Ghost's simple keyword stuffing, Fraud overhauls their entire CV. However, Bayesian Fusion restricts their final boost to just $+3.1$ points by reflecting the low signal strength of their actual GitHub profile. Note that there is still partial success since their $\Delta_{\mathrm{honest\_fusion}}$ score is $-2.4$, Bayesian Fusion still penalises them for the false claims.

\subsubsection{Documented Failure Modes}
\textbf{Smart Squatter.} The biggest failure mode is the Smart Squatter, where the candidate bypasses the Integrity Engine by finding a highly-scored GitHub portfolio that passes the fuzzy name-match. While rare in the real world, it remains a possible attack vector.

\textbf{Gamer.} Another impractical failure is the Gamer, who manipulates public-repository metrics to trick the system into thinking they were recently active ($\Delta_{\mathrm{fusion}} = +23.1$). This highlights a critical flaw in how MERIT calculates temporal skill decay.


\subsubsection{Study 06 takeaway}
The integrity engine is able to detect and penalise most forms of realistic adversarial attacks, such as keyword stuffing, identity squatting, and repository gaming.
Note that the majority of failure modes are edge-cases that take significant effort to achieve, and are not expected to be a major concern in the real world.
For example, let us consider the Smart Squatter scenario. For a given role, let us assume the malicious candidate wishes to squat a GitHub portfolio. First,
they would need to find portfolios that would score high on MERIT, which is a somewhat feasible task if they are willing to invest the time required to do this.
From here, the actor would need to find a name that is close enough to bypass the fuzzy name-match, this is likely not a guaranteed success due to the rarity of
a candidate finding a strong portfolio that scores high on MERIT for a specific role that shares the same name as the individual. This effort would technically work, 
but in the real world, where a candidate applies to a vast array of jobs with different metrics, this would be impractical. This effort would have to be repeated for
each role family that they apply to. In our example, we consider GitHub portfolios, but this identity squatting can be applied to other data sources, such as LinkedIn as well.

\subsection{Studies 02 \& 03: Scalability \& Performance Benchmarks}
\label{subsec:studies-02-03-scalability}

Studies~02 and~03 sit outside of the research questions in Table~\ref{tab:eval-literature-traceability}. They 
assess deployment feasibility (latency and memory at batch scale) rather than whether MERIT closes the extensibility, multi-source, 
or transparency gaps from Section~\ref{sec:relevance-to-merit}.
We benchmark only against our in-repository engine variants.


The runtime and spacetime studies are performed for cohorts of size $N \in \{10, 50, 100, 200, 500\}$ against the two baselines (Traditional and Modern AI),
as well as three variants of MERIT (CV-Only, Full without explainability, Full with explainability). In Study~02, we initialise each
engine once and then measure the wall-clock time for each engine under each cohort size. For study~03, we separate the initialisation 
cost from the incremental batch overhead per candidate. We run these benchmarks on Windows~11 (AMD Ryzen~7 7800X3D, Python~3.11.9).
Scoring is a purely CPU-bound operation, we do not leverage GPU acceleration or other hardware accelerations. We also do not
account for costs such as PDF ingestion, API calls and network latency from the interactions with our Supabase database.
The results are presented in Table~\ref{tab:runtime-results} and~\ref{tab:study03-memory}, 
for a graphical representation, see Appendix~\ref{appendix:fairness-scalability-charts}.

\begin{table}[htbp]
    \centering
    \caption[Study~02: Batch ranking latency]{Study~02: batch ranking latency (seconds) vs.\ cohort size $N$.}
    \label{tab:runtime-results}
    \vspace{1.5ex}
    \small
    \setlength{\tabcolsep}{6pt}
    \begin{tabular}{r c c c c c}
        \toprule
        \textbf{$N$} & \textbf{Trad. ATS} & \textbf{Mod. AI ATS} & \textbf{MERIT CV} & \textbf{MERIT Full} & \textbf{MERIT Expl.} \\ \midrule
        10  & 0.003  & 4.33   & 0.090 & 0.092 & 0.513 \\
        50  & 0.012  & 21.14  & 0.435 & 0.465 & 2.588 \\
        100 & 0.026  & 42.18  & 0.890 & 0.963 & 5.169 \\
        200 & 0.040  & 85.03  & 1.789 & 1.862 & 10.48 \\
        500 & 0.088  & 212.62 & 4.409 & 4.842 & 26.06 \\ \bottomrule
    \end{tabular}
\end{table}

\begin{table}[htbp]
    \centering
    \caption[Study~03: Engine memory]{Study~03: engine memory (MB). \textit{init.}\,=\,initialisation (static load at startup); remaining rows ($N$): incremental batch overhead after \texttt{gc.collect()} (not cumulative with static load).}
    \label{tab:study03-memory}
    \vspace{1.5ex}
    \small
    \setlength{\tabcolsep}{6pt}
    \begin{tabular}{r c c c c c}
        \toprule
        \textbf{$N$} & \textbf{Trad. ATS} & \textbf{Mod. AI ATS} & \textbf{MERIT CV} & \textbf{MERIT Full} & \textbf{MERIT Expl.} \\ \midrule
        \textit{init.} & 0.10 & 398.0 & 367.1 & 367.5 & 367.7 \\ \addlinespace
        10  & 0.02 & 113.0 & 54.4 & 54.5 & 53.8 \\
        50  & 0.03 & 116.7 & 56.3 & 56.7 & 56.7 \\
        100 & 0.05 & 118.5 & 58.8 & 59.1 & 59.4 \\
        200 & 0.06 & 119.7 & 63.3 & 63.5 & 65.0 \\
        500 & 0.12 & 119.1 & 73.9 & 74.6 & 76.5 \\ \bottomrule
    \end{tabular}
\end{table}

At $N = 500$, we see that MERIT Full is able to score the batch in under $4.4$ seconds, compared to $212.6$ seconds 
for the Modern AI ATS. Note the introduction of explainability, specifically Shapley values results in a runtime of 
$26.1$ seconds, leading to a $5.93$x slowdown compared to MERIT without explainability. There are $(2^3 = 8)$ possible 
coalitions of sources, and this will scale exponentially with the number of sources, not cohort size. Multi-source
fusion adds negligible latency and RSS overhead compared to CV only ($+0.43$ seconds and under $1$ MB at $N = 500$).
The majority of static load can be attributed to the SBERT initialisation (398.0 MB for Modern AI ATS, 367.1 MB for MERIT).
MERIT scales linearly with $N$ ($54.4$ MB to $73.9$ MB) because the harness retains per-candidate payloads, however,
this is still relatively modest for desktop batch screening.

\section{Threats to Validity}
\label{sec:threats-to-validity}

In previous sections of this chapter, we address our research questions (RQ1--RQ3) 
(Section~\ref{subsec:eval-design-rqs})
with twelve targeted studies. The other three studies are used to evaluate the fairness, 
security and computational scalability of MERIT. Like any other evaluation, we need to consider that these
results are subject to threats that can detract from any of the claims that we make. We list these threats
according to the works of Wohlin et al.~\cite{wohlin_threats}. More specifically, we have the following forms of
validity to which we need to consider threats: construct, internal, external and conclusion validity.
In this section, we state what each metric can, and cannot, support. 

\subsection{Construct Validity}
\label{subsec:construct-validity}

The first form of validity is that of construct validity, which refers to the extent
to which the metrics and instruments we use actually measure the research questions we claim to address.

\textbf{Synthetic ground truth.}
Some of our studies, specifically Studies~01, and 07--10, use synthetic ground truth data in the form of
JSON profiles and CV PDFs. Due to the nature of this data being created by the author, it is subject to
confirmation bias, the evaluation data we synthesise may unintentionally be authored to exhibit the
behaviours that we claim to test for, essentially cherry-picking the data to support our claims.
We aim to reduce the effect of confirmation bias by deliberately including cases that are designed to
fail for our system to provide a more robust evaluation. One study worth pointing out is the Signal
Dissonance study (Study~10), where we purposely create contradictory multi-source evidence to 
contradict the CV. None of our synthetic data is externally validated, and we do not use
real-world data due to GDPR and PII constraints. The controlled personas and synthetic data
were, therefore, the most ethical way to stress-test our system without the risk of exposing real candidates,
while maintaining reproducibility, since we include the data in the GitHub repository.

\textbf{Proxy metrics and the gap between ingestion and scoring.}
Throughout our studies, we report metrics such as parser F$_1$, Spearman $\rho$, shortlist rates, SUS, and more.
These metrics are not hire outcomes, nor interview pass rates or offer acceptance rates. Because of this, we
cannot claim that MERIT ``improves hiring outcomes''. A high Spearman $\rho$ or narrow Blind Mode prestige gap
cannot be interpreted as MERIT improving these decisions, however, we can still observe the nuanced behaviour
of our system and draw specific conclusions about the behaviour of our system. Due to the small sample size,
we cannot generalise the results to a larger population.

Another thread is that our Human-Machine Alignment studies (Study~07--10) use pre-parsed JSON data that assumes
perfect ingestion of data. In our RQ1 studies, we show that ingestion is not perfect, even under authored synthetic data,
therefore end-to-end MERIT behavior may differ from what we observe in these studies as they do not account for parser noise.

\textbf{Study-specific construct threats (01, 04--05, 06--11, 13--15).}
Study~01 and 06--10 intentionally use sparse LinkedIn personas to stress test the
contrast between a CV's claimed evidence and GitHub's demonstrated evidence. LinkedIn as a data source is considered
supplementary to a CV profile to account for parser error. However, since we assume perfect ingestion of data, we do not
need LinkedIn to fill in the gaps, therefore we can still observe the behaviour of our system without the need for LinkedIn
(Section~\ref{subsec:study-01-comparative-ablation}, 
Section~\ref{sec:human-machine-alignment}).

Study~04's skill recall does not reflect parser failure alone, it reflects the dictionary scope of the Devicon dataset
(Section~\ref{subsec:study-04}).
Study~05 bounds how much we can trust JD-driven metric configuration when JD extraction is imperfect (Section~\ref{subsec:study-05}).
Study~13 employs a single global corpus for all TF-IDF calculations, meaning that the results suffer
from cross-role IDF interference. The study also suffers from uneven distribution of scores from the participants who
help in the generation of the ground truth. Since TF-IDF aims for a fairly even distribution of scores, we depress
the $\rho$ because of these decisions (Section~\ref{subsubsec:study-13-limitations}).

The SUS scores generated from Study~14 indicate prototype feasibility, they are 
not product-market validation (Section~\ref{subsec:hci-trial-protocol}).
Blind Mode standardises layout, redacts names, employers and institutions, but it keeps tier labels due to its
usage in the scoring pipeline. Because of this, Study~15 does not guarantee demographic fairness.
We do not mask engine scores in Blind Mode, so evaluators may have anchored on match percentages for shortlist decisions (Section~\ref{subsec:bias-mitigation-efficacy}).
Blind Mode must not be interpreted as a fairness certificate.

\subsection{Internal Validity}
\label{subsec:internal-validity}

Internal validity concerns whether the observed results are caused by the variables 
under study, rather than by confounding factors.

\textbf{In-repository baselines and benchmarks.}
The baselines in our studies (Traditional ATS and Modern AI ATS) are in-repository implementations that aim to
proxy the keyword-matching and AI-assisted ATS categories mentioned earlier in Section~\ref{sec:automated-screening}.
They are not commercial products, nor do they replicate the full complexity or sophistication of the real-world ATS systems.
The keyword baseline may produce lower scores than the semantic baseline in real-world ATS systems.
The Modern AI baseline uses SBERT to compute cosine similarities between CV embeddings and JD embeddings, real world
ATS systems may employ more sophisticated models which would likely produce better results.
Reported baseline statistics for runtime and spacetime compare the algorithm families on the author's workstation.
Our baselines do not account for factors such as API latency, PDF ingestion, and network I/O.
Studies~02--03 only record single runs on a workstation without confidence intervals (Section~\ref{subsec:studies-02-03-scalability}).
Scores are likely to vary from run to run, however, the general pattern will remain stable across 
configurations under identical hardware and software.


\textbf{Protocol and scenario design.}
As noted under construct validity, our engines use pre-parsed JSON for evaluation, here we note
that the rank shifts do not account for ingestion errors, they do not validate the full PDF upload pipeline.
Also, Study~14 always runs the first stage of the pipeline before the second, so the reversals and confidence
gains might be influenced by order effects, rather than just explainability.
Study~06 clones one honest baseline and perturbs it with eight different scenarios. This cleanly isolates
each perturbation, but it does not tell us how often such attacks appear in real candidate pools (Section~\ref{subsec:study-06-adversarial}).
Studies~01 and~10 intentionally include engineered success and failure narratives to aid explanation 
(Sections~\ref{subsec:study-01-comparative-ablation}, 
\ref{subsec:study-10-spearman-dissonance}).

\subsection{External Validity}
\label{subsec:external-validity}

This form of validity revolves around determining the extent to which these findings can be generalised to the wider population,
beyond the conditions of these studies.

\textbf{Sample size and controlled scope.}

The cohort sizes for Studies~07--10 are small, more specifically, 10 candidates per cohort. Due to this
both the $\rho$ and $p$-values are volatile, and we cannot draw any conclusions about the generalisability 
of the results based on the $p$-values.
Studies~13--15 involve small panels of $N = 4$--$5$ participants, thus it cannot support a generalisation
at the population level.
In our evaluations, we often prioritise controlled ablation over field realism
since we want to isolate the behaviour of our system from the noise of the real world.
Synthetic stress tests and in-repository baselines help us isolate MERIT's
inner mechanisms, however, they do not replicate live ATS workflows, nor real hiring outcomes.

\textbf{Templates, evidence, and deployment.}
Study~04's multi-column layouts follow a specific template family, they do not reflect the
full diversity of multi-column layouts in the real world. MERIT's multi-source pillar rewards
public LinkedIn and GitHub signals in the majority of cases, therefore, we cannot generalise 
the results to the full population of candidates. Our integration of LinkedIn is not deployment-ready,
we need to benchmark it against an official LinkedIn API to be able to generalise the results.
Candidates without public evidence (Carl Corp, Section~\ref{subsec:study-01-comparative-ablation})
are often disadvantaged relative to juniors that have visible open-source histories.
Benchmarking doesn't account for the latency introduced by API calls, PDF ingestion, Network I/O, and database persistence.
The measurements we collect are not representative of production latency and memory usage.

\subsection{Conclusion Validity}
\label{subsec:conclusion-validity}

Conclusion validity addresses whether the conclusions we draw from our data are actually justified
by our choice of method and metrics.

\textbf{Spearman $\rho$ as the primary alignment metric.}

Spearman $\rho$ is a measurement of monotonic agreement against another ordering. It cannot measure
the calibration of match percentages, nor can it measure the magnitude of individual rank swaps.
Our engines may produce similar $\rho$ values, but they may produce significantly different score spreads.
We include the full rankings and scores for each engine in the project repository~\cite{merit_repo}, rather than
treating correlation coefficients alone as enough evidence. Studies~07--10 have recruiters rank only CV PDFs, 
while MERIT (Full) has access to GitHub and LinkedIn data. Spearman $\rho$ therefore reflects protocol
asymmetry along with engine quality (Section~\ref{subsec:spearman-method}).
The negative $\rho$ in Study~08 is due to engine misalignment under an intern JD, not sampling noise.
Our engine does not account for, nor claim seniority-aware hiring unless we introduce explicit policy rules
to account for this.

\textbf{Shapley verification and multiple studies.}
In Study~11, we verify efficiency, null-player, additivity, and determinism on a set of ten candidate profiles.
We intentionally leave out the axiom of symmetry since sources are heterogeneous by design (Section~\ref{subsec:study-11-shapley-verification}).
The negative GitHub attribution for Vince Vault is mathematically valid in our evaluation, however it
may be considered as a cognitive surprise to recruiters.

We run fifteen studies on overlapping personas and engines, which creates a risk of selectively emphasising supportive
results, even if we document failures explicitly.
The chapter should not be read as a meta-analysis with pooled effect sizes.

\subsection{Bounded Claims}
\label{subsec:validity-mitigations}

Despite the limitations we note in this section, the chapter still provides convergent directional evidence under
controlled data.
We therefore frame three bounded claims, one per research question.
Each claim states what the evaluation supports under our internal criteria (Table~\ref{tab:eval-literature-traceability}), and what it does not.


\textbf{RQ1 (Extensibility).}
Studies~04--05 allow us to claim that PDF ingestion is usable, but below the $\geq 80\%$ parsing target.
Skill recall is bounded by the scope of the Devicon dataset (38.6\%). Other metrics from this study are
influenced by layout sensitivity, especially for names on multi-column CVs (Section~\ref{subsec:study-04}).
Study~13 shows that TF-IDF metric priorities align with recruiter judgement directionally (aggregate $\rho = 0.25$), 
presenting exploratory evidence only (Section~\ref{subsubsec:study-13-limitations}).


\textbf{RQ2 (Multi-source fusion).}
Study~12 confirms that Beta--Bernoulli fusion favours demonstrated evidence over self-reported claims (Section~\ref{subsec:study-12-bayesian-conflict-resolution}).
Studies~01 and~06 show rank and score shifts that match design intent for hidden talent, fraud and
adversarial personas (Sections~\ref{subsec:study-01-comparative-ablation} and~\ref{subsec:study-06-adversarial}).
Studies~07--10 show that MERIT (Full) meets the directional $\rho \geq 0.7$ criterion in Studies~07 and~09,
however, this is not the case in Studies~08 or~10, and it does not lead every cohort column (Section~\ref{subsec:spearman-synthesis}).
We claim that multi-source fusion and integrity checks can improve alignment with expert ordering when
supplementary evidence agrees with the CV. We do not claim that MERIT (Full) replaces keyword or semantic
baselines in all hiring contexts, nor that it can perform well when candidates have no public evidence.
We cannot claim that MERIT works well when presented with batch-wide contradictory evidence, nor can we
claim MERIT can perform well when presented with nuances such as overqualification (Section~\ref{subsec:study-08-spearman-seniority}).

\textbf{RQ3 (Explainability and Transparency).}
Study~11 verifies the mathematical correctness of the Shapley values, and shows 
that they satisfy the efficiency, null-player, additivity, and determinism axioms on synthetic profiles (Section~\ref{subsec:study-11-shapley-verification}).
Negative GitHub attributions are also mathematically valid in that formulation.
Study~14 meets the SUS target in our pilot ($N=5$, mean SUS 81.5), and we observe decision reversals
throughout the study, presenting directional evidence for glass-box feasibility (Section~\ref{subsec:hci-trial-protocol}).
Study~15 shows directional compression of prestige-driven shortlist gaps in the UI (50~pp to 12.5~pp), but it is not
proof of demographic fairness due to the small-scale pilot (Section~\ref{subsec:bias-mitigation-efficacy}).
We claim that our transparency makes penalties and uplifts contestable for recruiters, but we do not 
claim that it removes automation bias since visible engine scores and tier labels can still anchor 
evaluators rather than eliminate bias (Section~\ref{subsec:construct-validity}).

\textbf{Supplementary: Feasibility and Security.}
Studies~02--03 both support deployment feasibility from a scalability and performance perspective under
three data sources. We can claim that MERIT's Shapley explainability layer is usable, but it is not cost-free at scale.
It is scalable with candidate size, but not with the number of data sources.
Study~06 shows that adversarial perturbations can be used to stress test the integrity of the system, and that
the system is able to detect and respond to the majority of these perturbations. We do not claim that MERIT
defends against all adversarial attacks, nor can we claim that it is a fair hiring system.

\textbf{Final Claims.}
MERIT is not fairness-certified and should not be considered as a fully autonomous hiring authority.
Future work requires counterbalanced HCI trials to account for order effects.
To make stronger claims, MERIT needs larger independent recruiter panels, repeated benchmarks with variance reporting
on various hardware and language configurations. MERIT would also need end-to-end PDF ingestion tests as well as 
field studies that can link MERIT rankings to interview, hire and performance outcomes.

MERIT is best described as an \textbf{auditable screening assistant}, where its purpose is
to help recruiters inspect evidence. Decisions should not be treated as the final decision, but rather as a hypothesis
to be tested and challenged by the recruiter.

